% ============================================================
% HSBI --- Quantitative Research Methods
% FINAL Mock Examination --- SET C (Lectures 1--12, comprehensive)
% Focus: ISLP Chapters 7, 8, 10, 13 (~60%) + Chapters 2--6 (~40%)
% Prof. Dr. Christoph Weisser --- Summer Semester 2026
% ------------------------------------------------------------
% SINGLE SOURCE with solutions toggle:
%   Exam only:
%     pdflatex -jobname=Final_Mock_Exam_C final_mock_exam_c.tex
%   With solutions:
%     pdflatex -jobname=Final_Mock_Exam_C_Solutions "\def\withsolutions{1}% ============================================================
% HSBI --- Quantitative Research Methods
% FINAL Mock Examination --- SET C (Lectures 1--12, comprehensive)
% Focus: ISLP Chapters 7, 8, 10, 13 (~60%) + Chapters 2--6 (~40%)
% Prof. Dr. Christoph Weisser --- Summer Semester 2026
% ------------------------------------------------------------
% SINGLE SOURCE with solutions toggle:
%   Exam only:
%     pdflatex -jobname=Final_Mock_Exam_C final_mock_exam_c.tex
%   With solutions:
%     pdflatex -jobname=Final_Mock_Exam_C_Solutions "\def\withsolutions{1}% ============================================================
% HSBI --- Quantitative Research Methods
% FINAL Mock Examination --- SET C (Lectures 1--12, comprehensive)
% Focus: ISLP Chapters 7, 8, 10, 13 (~60%) + Chapters 2--6 (~40%)
% Prof. Dr. Christoph Weisser --- Summer Semester 2026
% ------------------------------------------------------------
% SINGLE SOURCE with solutions toggle:
%   Exam only:
%     pdflatex -jobname=Final_Mock_Exam_C final_mock_exam_c.tex
%   With solutions:
%     pdflatex -jobname=Final_Mock_Exam_C_Solutions "\def\withsolutions{1}% ============================================================
% HSBI --- Quantitative Research Methods
% FINAL Mock Examination --- SET C (Lectures 1--12, comprehensive)
% Focus: ISLP Chapters 7, 8, 10, 13 (~60%) + Chapters 2--6 (~40%)
% Prof. Dr. Christoph Weisser --- Summer Semester 2026
% ------------------------------------------------------------
% SINGLE SOURCE with solutions toggle:
%   Exam only:
%     pdflatex -jobname=Final_Mock_Exam_C final_mock_exam_c.tex
%   With solutions:
%     pdflatex -jobname=Final_Mock_Exam_C_Solutions "\def\withsolutions{1}\input{final_mock_exam_c.tex}"
% ============================================================
\documentclass[11pt,a4paper]{article}
\usepackage[utf8]{inputenc}
\usepackage[T1]{fontenc}
\usepackage[english]{babel}
\usepackage[margin=2.2cm]{geometry}
\usepackage{amsmath,amssymb}
\usepackage{booktabs,array}
\usepackage{enumitem}
\usepackage{xcolor}
\usepackage{tcolorbox}
\tcbuselibrary{skins,breakable}

\setlength{\parindent}{0pt}

% Teal solution box (only shown when \withsolutions is defined)
\newtcolorbox{solutionbox}{enhanced,breakable,colback=teal!5,colframe=teal!55!black,
  boxrule=0pt,leftrule=2.5pt,arc=2pt,fonttitle=\bfseries,title=Solution,
  left=2mm,right=2mm,top=1mm,bottom=1.5mm}

\newcommand{\problem}[2]{\par\vspace{7mm}\noindent{\large\textbf{Problem #1}}%
  \hfill{\large\textbf{(#2 points)}}\par\nopagebreak\vspace{0.6mm}%
  \noindent\rule{\textwidth}{0.6pt}\par\nopagebreak\vspace{2.5mm}}
\newcommand{\pp}[1]{\textbf{[#1~P]}\enspace}

\begin{document}

% ------------------------------------------------------------
% Header block
% ------------------------------------------------------------
\begin{center}
  {\Large\textbf{HSBI --- Bielefeld University of Applied Sciences and Arts}}\\[1.2mm]
  {\large Quantitative Research Methods}\\[1mm]
  {\normalsize Summer Semester 2026 \quad$\cdot$\quad Examiner: Prof.\ Dr.\ Christoph Weisser}\\[4.5mm]
  {\LARGE\textbf{Final Mock Examination --- Set C (Lectures 1--12)}}\\[1.2mm]
  {\normalsize Comprehensive coverage of ISLP Chapters 2--8, 10 and 13,\\
   with emphasis on Chapters 7 (splines/GAMs), 8 (trees), 10 (deep learning) and 13 (multiple testing)}%
  \ifdefined\withsolutions\\[2.2mm]{\large\textcolor{teal!55!black}{\textbf{--- Version with detailed solutions ---}}}\fi
\end{center}

\vspace{2.5mm}
\begin{center}
\begin{tabular}{@{}ll@{}}
  \textbf{Duration:} & \textbf{120 minutes}\\
  \textbf{Total credit:} & \textbf{120 points}\\
  \textbf{Allowed aids:} & non-programmable calculator; one handwritten A4 sheet of notes (both sides)\\
\end{tabular}
\end{center}

\vspace{3mm}
Name: \rule[-1mm]{0.38\textwidth}{0.4pt}\hfill Matriculation no.: \rule[-1mm]{0.24\textwidth}{0.4pt}

\vspace{4.5mm}
\begin{center}
\renewcommand{\arraystretch}{1.35}
\begin{tabular}{|l|c|c|c|c|c|c|c|c|}
\hline
\textbf{Problem} & \textbf{P1} & \textbf{P2} & \textbf{P3} & \textbf{P4} & \textbf{P5} & \textbf{P6} & \textbf{P7} & \textbf{Total}\\\hline
Maximum points & 16 & 8 & 20 & 16 & 20 & 20 & 20 & \textbf{120}\\\hline
Achieved points & \hspace{8mm} & \hspace{8mm} & \hspace{8mm} & \hspace{8mm} & \hspace{8mm} & \hspace{8mm} & \hspace{8mm} & \hspace{10mm}\\\hline
\end{tabular}
\end{center}

\vspace{2.5mm}
\textbf{Instructions}
\begin{itemize}[nosep,leftmargin=5.5mm]
  \item Justify all answers. Numerical results without a derivation or a brief reasoning receive no credit.
  \item Round final numerical results to 3 decimal places unless stated otherwise. Small deviations caused by carrying rounded intermediate results are accepted.
  \item All numerical constants you need (powers of $e$, logarithms, $(0.95)^{40}$, $t$-quantiles) are provided in the problems.
  \item The number of points per problem roughly equals the recommended working time in minutes.
\end{itemize}
\vspace{1mm}\noindent\rule{\textwidth}{0.6pt}

% ============================================================
\problem{1: Moving beyond linearity --- splines}{16}
{\small\itshape Lecture slides: Chapter 7 --- Moving Beyond Linearity (Lecture 9).}\par\vspace{1mm}

\begin{enumerate}[label=(\alph*),leftmargin=7mm,itemsep=2.5mm]
  \item \pp{4} A \emph{cubic spline} is a piecewise cubic polynomial that is continuous and has
    continuous first and second derivatives at each knot. Derive the number of degrees of freedom
    (free parameters, counting the intercept) of a cubic spline with $K$ interior knots by
    counting the coefficients of the polynomial pieces and subtracting the constraints imposed
    at the knots. Evaluate your formula for $K=3$ interior knots.
  \item \pp{4} How many degrees of freedom does a \emph{natural} cubic spline with the same
    $K=3$ knots use? State the additional constraint a natural spline imposes, say in which
    region of the predictor space it takes effect, and explain why it is useful in practice.
  \item \pp{4} A \emph{smoothing spline} is the function $g$ that minimises
    \[
      \sum_{i=1}^{n}\bigl(y_i-g(x_i)\bigr)^2+\lambda\int g''(t)^2\,dt .
    \]
    Explain the role of each of the two terms, and describe the behaviour of the solution
    $\hat{g}$ in the two limiting cases $\lambda\to 0$ and $\lambda\to\infty$.
  \item \pp{4} Write down the \emph{truncated-power basis} representation of a cubic spline with
    a single knot at $\xi$, define the truncated-power basis function explicitly, and explain in
    one sentence why adding this basis function does not destroy the smoothness of the fit at
    $\xi$.
\end{enumerate}

\ifdefined\withsolutions
\begin{solutionbox}
\textbf{(a) Degrees of freedom of a cubic spline.}
$K$ interior knots split the range of $X$ into $K+1$ intervals, each carrying a cubic
polynomial with $4$ coefficients: $4(K+1)=4K+4$ parameters. At each knot, continuity of
$g$, $g'$ and $g''$ imposes $3$ constraints, i.e.\ $3K$ constraints in total. Hence
\[
  \mathrm{df}=4(K+1)-3K=K+4 .
\]
For $K=3$: $\mathrm{df}=3+4=\mathbf{7}$ free parameters (equivalently: intercept, $x$, $x^2$,
$x^3$ plus one truncated-power basis function per knot $=4+K$).

\medskip
\textbf{(b) Natural cubic spline.}
A natural spline additionally requires the function to be \emph{linear beyond the boundary
knots}, i.e.\ in the two outer regions $X<\xi_1$ and $X>\xi_K$. Forcing the quadratic and cubic
terms to vanish there removes $2$ parameters per boundary region, i.e.\ $4$ in total:
\[
  \mathrm{df}=(K+4)-4=K,\qquad K=3:\ \mathrm{df}=\mathbf{3}.
\]
Usefulness: polynomials and unconstrained splines are notoriously erratic near the boundaries,
where data are sparse; the linearity constraint stabilises the fit there
(\emph{lower variance} of predictions at and beyond the boundary), at the price of a small
amount of additional bias.

\medskip
\textbf{(c) Smoothing spline.}
The first term (RSS) rewards fidelity to the data; the second term penalises roughness, since
$g''$ measures the curvature (wiggliness) of $g$, and $\lambda\ge 0$ is the tuning parameter
that trades the two off.
\emph{$\lambda\to 0$:} the penalty disappears; $\hat{g}$ can be made arbitrarily wiggly and
\emph{interpolates} the training observations exactly (zero training RSS, severe overfitting;
effective df $\to n$).
\emph{$\lambda\to\infty$:} any curvature is infinitely expensive, so $g''\equiv 0$ is enforced;
$\hat{g}$ becomes a straight line and equals the \emph{least-squares regression line}
(effective df $\to 2$). Between the extremes, $\lambda$ (chosen e.g.\ by LOOCV) controls the
effective degrees of freedom, which decrease smoothly from $n$ to $2$.

\medskip
\textbf{(d) Truncated-power basis with one knot $\xi$.}
\[
  f(x)=\beta_0+\beta_1x+\beta_2x^2+\beta_3x^3+\beta_4\,(x-\xi)^3_+,
  \qquad
  (x-\xi)^3_+=\begin{cases}(x-\xi)^3, & x>\xi\\[0.5mm] 0, & x\le\xi\end{cases}
\]
The function $(x-\xi)^3_+$ and its first and second derivatives are all $0$ at $x=\xi$, so
adding it changes only the third derivative at the knot; the fit remains continuous with
continuous first and second derivatives --- exactly the defining smoothness of a cubic spline.
\end{solutionbox}
\fi

% ============================================================
\problem{2: Generalised additive models}{8}
{\small\itshape Lecture slides: Chapter 7 --- Moving Beyond Linearity (Lecture 9).}\par\vspace{1mm}

\begin{enumerate}[label=(\alph*),leftmargin=7mm,itemsep=2.5mm]
  \item \pp{3} Write down the general form of a \emph{generalised additive model} (GAM) for a
    quantitative response $Y$ and predictors $X_1,\dots,X_p$, and explain in one sentence how it
    generalises multiple linear regression.
  \item \pp{3} State one important \emph{advantage} of GAMs concerning the interpretation of the
    fitted functions $f_j$, and the main structural \emph{limitation} of the additive form ---
    including how this limitation can be repaired.
  \item \pp{2} Describe in one or two sentences how the \emph{backfitting} algorithm fits a GAM.
\end{enumerate}

\ifdefined\withsolutions
\begin{solutionbox}
\textbf{(a) Model form.}
\[
  Y=\beta_0+f_1(X_1)+f_2(X_2)+\dots+f_p(X_p)+\varepsilon .
\]
Multiple linear regression is the special case $f_j(X_j)=\beta_jX_j$; a GAM replaces each linear
component by a smooth non-linear function $f_j$ (e.g.\ a natural or smoothing spline, or local
regression), while keeping the \emph{additive} structure.

\medskip
\textbf{(b) Advantage and limitation.}
\emph{Advantage (interpretability):} because the model is additive, the contribution of each
$X_j$ to $Y$ is captured by its own function $f_j$, which can be examined and plotted
individually while the other variables are held fixed --- non-linear effects remain as easy to
read as in a coefficient table, one panel per predictor.
\emph{Limitation:} additivity excludes \emph{interactions}: the effect of $X_j$ cannot depend
on the value of $X_k$. Repair: add interaction terms manually, e.g.\ a product term
$X_j\times X_k$ or a low-dimensional smooth surface $f_{jk}(X_j,X_k)$ as an extra component.

\medskip
\textbf{(c) Backfitting.}
Cycle through the predictors and repeatedly re-fit each $f_j$ (by a smoother, e.g.\ a spline)
to the \emph{partial residuals} $r_i=y_i-\beta_0-\sum_{k\ne j}f_k(x_{ik})$, holding all other
fitted functions fixed; iterate until the $f_j$ stop changing.
\end{solutionbox}
\fi

% ============================================================
\problem{3: Regression and classification trees by hand}{20}
{\small\itshape Lecture slides: Chapter 8 --- Tree-Based Methods (Lecture 10).}\par\vspace{1mm}

A logistics company records for $n=6$ regional depots the number of delivery vans in operation
$x$ and the number of parcels delivered that day $y$ (in 1{,}000 parcels):

\begin{center}
\begin{tabular}{lcccccc}
\toprule
Depot $i$ & 1 & 2 & 3 & 4 & 5 & 6\\
\midrule
Vans $x_i$ & 1 & 2 & 3 & 4 & 5 & 6\\
Parcels $y_i$ & 3 & 5 & 4 & 9 & 11 & 10\\
\bottomrule
\end{tabular}
\end{center}

A regression tree with a \emph{single} split at cutpoint $s$ partitions the data into
$R_1=\{i: x_i<s\}$ and $R_2=\{i: x_i\ge s\}$ and predicts the region mean $\hat{y}_{R_m}$ in
each region.

\begin{enumerate}[label=(\alph*),leftmargin=7mm,itemsep=2.5mm]
  \item \pp{8} For each of the two candidate cutpoints $s=2.5$ and $s=3.5$, determine the two
    regions, the region means and the resulting total residual sum of squares
    $\operatorname{RSS}(s)=\sum_{m=1}^{2}\sum_{i\in R_m}(y_i-\hat{y}_{R_m})^2$.
    Which cutpoint does recursive binary splitting choose, and why?
  \item \pp{2} Using the chosen split, predict the number of parcels for a new depot that will
    operate $x=5$ vans.
  \item \pp{6} A \emph{classification} tree for a different task has a terminal node containing
    $9$ training observations from three product categories: $5$ ``Electronics'', $3$
    ``Clothing'' and $1$ ``Books''. Compute the class proportions, the Gini index
    $G=\sum_{k}\hat{p}_{k}(1-\hat{p}_{k})$ and the entropy $D=-\sum_{k}\hat{p}_{k}\ln\hat{p}_{k}$
    of this node \emph{(use $\ln(5/9)=-0.5878$, $\ln(1/3)=-1.0986$ and $\ln(1/9)=-2.1972$)}.
    Which class does the node predict?
  \item \pp{4} In \emph{cost-complexity pruning} one minimises, for a subtree
    $T\subseteq T_0$ of the large tree $T_0$,
    \[
      \sum_{m=1}^{|T|}\ \sum_{i:\,x_i\in R_m}\bigl(y_i-\hat{y}_{R_m}\bigr)^2+\alpha\,|T| .
    \]
    Explain the role of the tuning parameter $\alpha$: what happens for $\alpha=0$, what happens
    as $\alpha$ increases, and how is $\alpha$ chosen in practice?
\end{enumerate}

\ifdefined\withsolutions
\begin{solutionbox}
\textbf{(a) Evaluating the two cutpoints.}
Root node for reference: $\bar{y}=42/6=7$.

\emph{Cutpoint $s=2.5$:} $R_1=\{1,2\}$ with $y$-values $\{3,5\}$, mean $\hat{y}_{R_1}=4$;
$R_2=\{3,4,5,6\}$ with $y$-values $\{4,9,11,10\}$, mean $\hat{y}_{R_2}=34/4=8.5$.
\begin{align*}
  \operatorname{RSS}_{R_1}&=(3-4)^2+(5-4)^2=2,\\
  \operatorname{RSS}_{R_2}&=(4-8.5)^2+(9-8.5)^2+(11-8.5)^2+(10-8.5)^2=20.25+0.25+6.25+2.25=29,
\end{align*}
so $\operatorname{RSS}(2.5)=2+29=\mathbf{31}$.

\emph{Cutpoint $s=3.5$:} $R_1=\{1,2,3\}$ with $y$-values $\{3,5,4\}$, mean
$\hat{y}_{R_1}=12/3=4$; $R_2=\{4,5,6\}$ with $y$-values $\{9,11,10\}$, mean
$\hat{y}_{R_2}=30/3=10$.
\[
  \operatorname{RSS}_{R_1}=(3-4)^2+(5-4)^2+(4-4)^2=2,\qquad
  \operatorname{RSS}_{R_2}=(9-10)^2+(11-10)^2+(10-10)^2=2,
\]
so $\operatorname{RSS}(3.5)=2+2=\mathbf{4}$.

Since $4<31$, recursive binary splitting is \emph{greedy} and picks the cutpoint with the
smallest RSS: $\mathbf{s=3.5}$. (It cleanly separates the three low-volume depots, mean $4$,
from the three high-volume depots, mean $10$; in fact $s=3.5$ is the best of all five possible
cutpoints.)

\medskip
\textbf{(b) Prediction for $x=5$.}
$x=5\ge 3.5$, so the new depot falls into $R_2$ and the tree predicts
$\hat{y}=\hat{y}_{R_2}=\mathbf{10}$ (i.e.\ $10{,}000$ parcels).

\medskip
\textbf{(c) Node impurity (three classes).}
$\hat{p}_{\text{E}}=5/9$, $\hat{p}_{\text{C}}=3/9=1/3$, $\hat{p}_{\text{B}}=1/9$.
\[
  G=\tfrac{5}{9}\cdot\tfrac{4}{9}+\tfrac{3}{9}\cdot\tfrac{6}{9}+\tfrac{1}{9}\cdot\tfrac{8}{9}
   =\tfrac{20+18+8}{81}=\tfrac{46}{81}=\mathbf{0.568}.
\]
\[
  D=-\Bigl(\tfrac{5}{9}\ln\tfrac{5}{9}+\tfrac{1}{3}\ln\tfrac{1}{3}+\tfrac{1}{9}\ln\tfrac{1}{9}\Bigr)
   =-\bigl(\tfrac{5}{9}(-0.5878)+\tfrac{1}{3}(-1.0986)+\tfrac{1}{9}(-2.1972)\bigr)
\]
\[
   =0.3266+0.3662+0.2441=\mathbf{0.937}.
\]
The node predicts the majority class, \textbf{Electronics} (with estimated probability
$5/9\approx 0.556$). Both $G$ and $D$ would be $0$ for a pure node; here the node is quite
impure because no class dominates strongly.

\medskip
\textbf{(d) Role of $\alpha$.}
$\alpha$ is the price per terminal node: it trades off the fit of the subtree (first term)
against its complexity $|T|$.
For $\alpha=0$ the criterion equals the training RSS and the full tree $T_0$ is optimal.
As $\alpha$ increases, additional leaves must ``pay for themselves''; branches whose RSS
reduction is smaller than $\alpha$ are pruned, so the optimal subtree becomes smaller ---
one obtains a \emph{nested sequence} of subtrees as $\alpha$ grows (weakest-link pruning).
In practice $\alpha$ is selected by \emph{cross-validation} (choose the $\alpha$ with the
smallest CV error, then refit on all data), which balances bias against variance instead of
optimising training fit.
\end{solutionbox}
\fi

% ============================================================
\problem{4: Ensembles --- bagging, random forests, boosting}{16}
{\small\itshape Lecture slides: Chapter 8 --- Tree-Based Methods (Lecture 10).}\par\vspace{1mm}

\begin{enumerate}[label=(\alph*),leftmargin=7mm,itemsep=2.5mm]
  \item \pp{6} Let $\hat{f}_1(x),\dots,\hat{f}_B(x)$ be $B$ identically distributed (but not
    independent) tree predictions at a fixed point $x$, each with variance $\sigma^2$ and with
    pairwise correlation $\rho$. The variance of the bagged average is
    \[
      \operatorname{Var}\Bigl(\tfrac{1}{B}\sum_{b=1}^{B}\hat{f}_b(x)\Bigr)
      =\rho\,\sigma^{2}+\frac{1-\rho}{B}\,\sigma^{2}.
    \]
    (i) Explain briefly why this formula shows that bagging reduces variance, and which term
    limits the reduction. (ii) Evaluate the variance for $B=50$, $\rho=0.4$ and $\sigma^2=16$,
    and compare it with the variance of a single tree. (iii) What value does the variance
    approach as $B\to\infty$?
  \item \pp{4} How do \emph{random forests} modify bagging at each split, and what is the
    recommended subset size $m$ for \emph{regression} as a function of $p$ (evaluate for
    $p=30$)? Explain, with reference to the formula in (a), why this modification improves on
    bagging.
  \item \pp{4} Name the three tuning parameters of \emph{boosting} for regression trees and
    state in one sentence each what they control. What is the ``slow learning'' idea behind
    boosting?
  \item \pp{2} What does the \emph{out-of-bag} (OOB) error of a bagged model estimate, and why
    is it available essentially for free?
\end{enumerate}

\ifdefined\withsolutions
\begin{solutionbox}
\textbf{(a) Variance of the bagged average.}
(i) The second term $\frac{1-\rho}{B}\sigma^2$ is driven to $0$ by averaging over many trees
($B$ large), so averaging removes the ``independent part'' of the trees' variability ---
this is the variance reduction of bagging. The first term $\rho\sigma^2$ does \emph{not}
depend on $B$: because the bootstrap trees are grown on overlapping samples of the same data
(and share the same strong predictors), they are positively correlated, and this correlated
part of the variance cannot be averaged away. It limits the reduction.
(ii) For $B=50$, $\rho=0.4$, $\sigma^2=16$:
\[
  \operatorname{Var}=0.4\cdot 16+\frac{1-0.4}{50}\cdot 16
  =6.4+\frac{0.6\cdot 16}{50}=6.4+0.192=\mathbf{6.592},
\]
compared with $\sigma^2=16$ for a single tree --- a reduction of about $58.8\%$.
(iii) As $B\to\infty$ the variance approaches the floor $\rho\sigma^2=0.4\cdot 16=\mathbf{6.4}$.
(Increasing $B$ further does not overfit; it only stabilises the average.)

\medskip
\textbf{(b) Random forests.}
At each split, only a \emph{random subset} of $m$ of the $p$ predictors is considered as
splitting candidates; a fresh subset is drawn at every split. Recommended for regression:
$m\approx p/3$, so for $p=30$: $m=30/3=\mathbf{10}$.
Why it helps: with bagging, a dominant predictor is used in the top split of almost every tree,
so the trees look alike and $\rho$ is high. Excluding it from $\approx (p-m)/p$ of the split
searches \emph{decorrelates} the trees, i.e.\ lowers $\rho$ --- and by (a) this lowers the
irreducible floor $\rho\sigma^2$ of the ensemble variance, which mere averaging cannot touch.

\medskip
\textbf{(c) Boosting.}
\begin{itemize}[nosep,leftmargin=5mm]
  \item \emph{Number of trees $B$:} how many small trees are added sequentially; unlike
    bagging, boosting \emph{can} overfit if $B$ is too large, so $B$ is chosen by
    cross-validation.
  \item \emph{Shrinkage $\lambda$} (learning rate, e.g.\ $0.01$ or $0.001$): scales the
    contribution of each new tree and thus the speed of learning; smaller $\lambda$ usually
    needs larger $B$.
  \item \emph{Number of splits $d$ per tree} (interaction depth): the complexity of each tree
    and the interaction order of the model; often even $d=1$ (stumps, additive model) works
    well.
\end{itemize}
\emph{Slow learning:} each tree is fitted to the current \emph{residuals}, and only a small
fraction $\lambda$ of it is added to the model; the model thus improves slowly and exactly in
the regions where it still performs poorly, instead of fitting the data hard in one step ---
which tends to generalise better.

\medskip
\textbf{(d) OOB error.}
Each bootstrap sample leaves out on average about $1/3$ of the observations
(``out-of-bag''). Predicting each observation only with the trees that did \emph{not} use it,
and aggregating these predictions, yields the OOB error --- a valid estimate of the
\emph{test error} of the bagged model. It is free because it reuses the bootstrap side
product; no separate validation set or cross-validation loop is required.
\end{solutionbox}
\fi

% ============================================================
\problem{5: Neural networks by hand}{20}
{\small\itshape Lecture slides: Chapter 10 --- Deep Learning (Lecture 11).}\par\vspace{1mm}

Consider a tiny feed-forward network for a quantitative response with input
$x=(x_1,x_2)=(2,3)$, one hidden layer with two ReLU units ($g(z)=\max\{0,z\}$) and a linear
output unit:
\[
  A_1=g\bigl(-3+2\,x_1+1\,x_2\bigr),\qquad
  A_2=g\bigl(2-2\,x_1-1\,x_2\bigr),\qquad
  f(x)=-2+3\,A_1+2\,A_2 .
\]

\begin{enumerate}[label=(\alph*),leftmargin=7mm,itemsep=2.5mm]
  \item \pp{8} Perform the forward pass step by step: compute both pre-activations, both
    activations $A_1,A_2$ and the network output $f(x)$. Comment briefly on what the ReLU did
    to the second unit and why such non-linear activations are essential.
  \item \pp{6} A classification network with $K=3$ classes produces the logits (output-layer
    pre-activations) $z=(z_1,z_2,z_3)=(2.0,\ 0.0,\ -1.0)$ for one observation whose true class
    is class~2. Compute all three softmax probabilities and the cross-entropy loss
    $-\ln \hat{p}_{\text{true}}$.
    \emph{(Use $e^{2}=7.389$, $e^{0}=1.000$, $e^{-1}=0.368$ and $\ln 8.757=2.170$.)}
  \item \pp{4} Count the parameters of a fully connected network with $p=6$ inputs, one hidden
    layer with $k=5$ units and an output layer with $K=3$ units, where every hidden and output
    unit has a bias. Show the count separately for the two weight layers.
  \item \pp{2} A convolutional layer receives a $64\times 64$ image and uses filters of size
    $F=3$ with padding $P=1$ and stride $S=2$. Compute the spatial size of the output feature
    map from $\bigl\lfloor(W-F+2P)/S\bigr\rfloor+1$ (note the stride is $2$, so the division is
    rounded down).
\end{enumerate}

\ifdefined\withsolutions
\begin{solutionbox}
\textbf{(a) Forward pass.}
Pre-activations:
\[
  z_1=-3+2\cdot 2+1\cdot 3=-3+4+3=4,\qquad
  z_2=2-2\cdot 2-1\cdot 3=2-4-3=-5 .
\]
ReLU activations:
\[
  A_1=g(4)=\max\{0,4\}=4,\qquad A_2=g(-5)=\max\{0,-5\}=0 .
\]
Output:
\[
  f(x)=-2+3\cdot 4+2\cdot 0=-2+12+0=\mathbf{10}.
\]
The ReLU \emph{switched the second unit off}: its negative pre-activation is clipped to $0$,
so unit 2 contributes nothing for this input (but would for other inputs). This thresholding
is the source of non-linearity: without a non-linear $g$, the network would collapse into a
single linear function of $x$, no matter how many layers it has.

\medskip
\textbf{(b) Softmax and cross-entropy.}
Exponentials: $e^{2.0}=7.389$, $e^{0.0}=1.000$, $e^{-1.0}=0.368$; their sum is
$7.389+1.000+0.368=8.757$. Softmax probabilities:
\[
  \hat{p}_1=\frac{7.389}{8.757}=0.844,\qquad
  \hat{p}_2=\frac{1.000}{8.757}=0.114,\qquad
  \hat{p}_3=\frac{0.368}{8.757}=0.042,
\]
which sum to $1.000$. Cross-entropy loss for the true class $2$:
\[
  -\ln\hat{p}_2=-\ln\frac{1.000}{8.757}=-\bigl(0-\ln 8.757\bigr)=\ln 8.757=\mathbf{2.170}.
\]
(Note that the model's most probable class is class~1, but the loss is evaluated at the
\emph{true} class~2, which received only probability $0.114$ --- hence the sizeable loss.)

\medskip
\textbf{(c) Parameter count for $6\to 5\to 3$.}
Input $\to$ hidden: $6\cdot 5=30$ weights $+\,5$ biases $=35$.
Hidden $\to$ output: $5\cdot 3=15$ weights $+\,3$ biases $=18$.
Total: $35+18=\mathbf{53}$ parameters.

\medskip
\textbf{(d) Convolution output size.}
\[
  \Bigl\lfloor\frac{W-F+2P}{S}\Bigr\rfloor+1
  =\Bigl\lfloor\frac{64-3+2\cdot 1}{2}\Bigr\rfloor+1
  =\Bigl\lfloor\frac{63}{2}\Bigr\rfloor+1=\lfloor 31.5\rfloor+1=31+1=32,
\]
so the output feature map is $\mathbf{32\times 32}$ --- the stride $S=2$ roughly halves the
$64\times 64$ input (the floor discards the fractional part before adding $1$).
\end{solutionbox}
\fi

% ============================================================
\problem{6: Multiple testing}{20}
{\small\itshape Lecture slides: Chapter 13 --- Multiple Testing (Lecture 12).}\par\vspace{1mm}

\begin{enumerate}[label=(\alph*),leftmargin=7mm,itemsep=2.5mm]
  \item \pp{4} A lab performs $m=40$ \emph{independent} hypothesis tests, each at level
    $\alpha=0.05$, and suppose all $40$ null hypotheses are true. Define the family-wise error
    rate (FWER) and compute it for this situation \emph{(use $(0.95)^{40}=0.1285$)}. Interpret
    the result in one sentence.
  \item \pp{7} A research group screens $m=6$ candidate biomarkers; the ordered $p$-values are
    \begin{center}
    \begin{tabular}{lcccccc}
    \toprule
    $j$ & 1 & 2 & 3 & 4 & 5 & 6\\
    \midrule
    $p_{(j)}$ & 0.001 & 0.008 & 0.011 & 0.030 & 0.041 & 0.260\\
    \bottomrule
    \end{tabular}
    \end{center}
    Control the FWER at $\alpha=0.05$ (i) with \emph{Bonferroni} and (ii) with the
    \emph{Holm step-down} procedure. Give the relevant thresholds, state for each method
    exactly which hypotheses are rejected, and explain why Holm can never reject fewer
    hypotheses than Bonferroni.
  \item \pp{6} Apply the \emph{Benjamini--Hochberg} (BH) procedure with FDR level $q=0.05$ to
    the same six $p$-values: give the thresholds $jq/m$, determine which hypotheses are
    rejected, and state precisely what ``FDR control at $q=0.05$'' guarantees for the
    resulting set of rejections.
  \item \pp{3} A pharmaceutical company runs a \emph{confirmatory} phase-III trial that tests
    $m=5$ pre-specified endpoints; a single false-positive endpoint could lead to the approval
    of an ineffective drug. Should they control the FWER or the FDR? Justify your
    recommendation.
\end{enumerate}

\ifdefined\withsolutions
\begin{solutionbox}
\textbf{(a) FWER for $m=40$ independent tests.}
$\text{FWER}=\Pr(\text{at least one Type I error})=\Pr(V\ge 1)$, where $V$ is the number of
true nulls that are (falsely) rejected. With independence and all nulls true:
\[
  \text{FWER}=1-\Pr(\text{no false rejection})=1-(1-0.05)^{40}=1-(0.95)^{40}
  =1-0.1285=\mathbf{0.8715}.
\]
Interpretation: even though each single test is performed at the $5\%$ level, there is about an
$87\%$ chance of at least one false discovery among the $40$ tests --- testing many hypotheses
without adjustment almost guarantees spurious findings.

\medskip
\textbf{(b) Bonferroni and Holm at $\alpha=0.05$.}
\emph{Bonferroni:} reject $H_{(j)}$ iff $p_{(j)}\le\alpha/m=0.05/6=0.00833$.
$p_{(1)}=0.001\le 0.00833$ \checkmark, $p_{(2)}=0.008\le 0.00833$ \checkmark,
$p_{(3)}=0.011>0.00833$ $\times$.
$\Rightarrow$ Bonferroni rejects $H_{(1)}$ and $H_{(2)}$: \textbf{2 rejections}.

\emph{Holm:} compare $p_{(j)}$ with $\alpha/(m-j+1)$ step by step and stop at the first failure:
\begin{center}
\begin{tabular}{lcccccc}
\toprule
$j$ & 1 & 2 & 3 & 4 & 5 & 6\\
\midrule
$p_{(j)}$ & 0.001 & 0.008 & 0.011 & 0.030 & 0.041 & 0.260\\
$\alpha/(m-j+1)$ & 0.00833 & 0.0100 & 0.0125 & 0.0167 & 0.0250 & 0.0500\\
reject? & \checkmark & \checkmark & \checkmark & stop & --- & ---\\
\bottomrule
\end{tabular}
\end{center}
$0.001\le 0.00833$, $0.008\le 0.0100$, $0.011\le 0.0125$, but $0.030>0.0167$: stop.
$\Rightarrow$ Holm rejects $H_{(1)},H_{(2)},H_{(3)}$: \textbf{3 rejections}.

\emph{Why Holm $\supseteq$ Bonferroni:} Holm's thresholds $\alpha/(m-j+1)$ are $\ge\alpha/m$
for every $j$ (with equality only at $j=1$), so every $p$-value that passes Bonferroni also
passes Holm's less stringent later hurdles --- Holm is uniformly more powerful, yet still
controls the FWER at $\alpha$ without any independence assumption. Here Holm rejects one more
hypothesis ($H_{(3)}$) than Bonferroni.

\medskip
\textbf{(c) Benjamini--Hochberg at $q=0.05$.}
Thresholds $jq/m=j\cdot 0.05/6$:
\begin{center}
\begin{tabular}{lcccccc}
\toprule
$j$ & 1 & 2 & 3 & 4 & 5 & 6\\
\midrule
$p_{(j)}$ & 0.001 & 0.008 & 0.011 & 0.030 & 0.041 & 0.260\\
$jq/m$ & 0.00833 & 0.0167 & 0.0250 & 0.0333 & 0.0417 & 0.0500\\
$p_{(j)}\le jq/m$? & yes & yes & yes & yes & \textbf{yes} & no\\
\bottomrule
\end{tabular}
\end{center}
$L=$ largest $j$ with $p_{(j)}\le jq/m$: $j=5$ (since $0.041\le 0.0417$, while
$0.260>0.0500$). BH rejects \emph{all} $H_{(j)}$ with $j\le L$: $H_{(1)},\dots,H_{(5)}$:
\textbf{5 rejections}.
Guarantee: $\text{FDR}=E\bigl[V/\max(R,1)\bigr]\le q=0.05$ --- \emph{on average} at most
$5\%$ of the rejected hypotheses are false discoveries. It does \emph{not} guarantee that at
most $5\%$ of these particular $5$ rejections are false, nor does it bound the probability of
making any false rejection (that would be FWER control). Note BH is the most liberal of the
three here ($5$ rejections vs.\ $3$ for Holm and $2$ for Bonferroni).

\medskip
\textbf{(d) Confirmatory phase-III trial.}
Control the \textbf{FWER}. This is a confirmatory study with only $m=5$ pre-specified
endpoints, and the cost of even a single false positive is severe (approval of an ineffective
drug, patient harm, regulatory failure). FWER control (e.g.\ Holm at $\alpha=0.05$) bounds the
probability of \emph{any} false rejection at $5\%$, which is exactly the guarantee a regulator
demands. FDR control would only limit the \emph{expected proportion} of false endpoints and
would tolerate an occasional false positive --- acceptable in large exploratory screens, but
not when each individual claim must stand on its own.
\end{solutionbox}
\fi

% ============================================================
\problem{7: Cumulative essentials (Chapters 2--6)}{20}
{\small\itshape Lecture slides: Chapters 2--6 (Lectures 1--8).}\par\vspace{1mm}

\begin{enumerate}[label=(\alph*),leftmargin=7mm,itemsep=2.5mm]
  \item \pp{4} The \emph{bootstrap} draws a resample of size $n$ with replacement from the $n$
    observations. (i) Show that the probability that a given observation is \emph{not} selected
    into a particular bootstrap sample equals $(1-1/n)^n$, and give its limit as $n\to\infty$
    \emph{(use $e^{-1}=0.368$)}. (ii) Hence what fraction of the original observations appears
    (at least once) in a typical bootstrap sample as $n\to\infty$? (iii) Compute this inclusion
    probability exactly for $n=5$ \emph{(use $0.8^{5}=0.328$)}.
  \item \pp{4} At a fixed test point $x_0$ the expected test error of an estimate $\hat f$
    decomposes as
    $E\bigl[(y_0-\hat f(x_0))^2\bigr]
      =\operatorname{Var}(\hat f(x_0))+\bigl[\operatorname{Bias}(\hat f(x_0))\bigr]^2
       +\operatorname{Var}(\varepsilon)$.
    Two candidate methods give, at $x_0$:
    \begin{center}
    \begin{tabular}{lccc}
    \toprule
    Method & $\operatorname{Var}(\hat f)$ & $\operatorname{Bias}(\hat f)$ & $\operatorname{Var}(\varepsilon)$\\
    \midrule
    (1) flexible spline & 3.5 & $-1.0$ & 2.0\\
    (2) linear model & 0.5 & $-2.5$ & 2.0\\
    \bottomrule
    \end{tabular}
    \end{center}
    Compute the expected test MSE of each method, say which you would deploy at $x_0$, and state
    the lowest test MSE \emph{any} method could achieve here.
  \item \pp{3} Write down the \emph{elastic net} penalty and explain what it inherits from the
    lasso and what it inherits from ridge regression. In which situation does the elastic net
    clearly outperform the lasso?
  \item \pp{3} Two predictors in a regression have correlation $0.95$. Explain what such
    \emph{collinearity} does to the least-squares coefficient estimates. The variance inflation
    factor is $\mathrm{VIF}_j=1/(1-R_j^2)$, where $R_j^2$ is from regressing $X_j$ on the other
    predictors; compute and interpret $\mathrm{VIF}_j$ for $R_j^2=0.9$ \emph{(use
    $\sqrt{10}=3.162$)}.
  \item \pp{6} Match each scenario to exactly one method from
    \{\emph{linear regression}, \emph{lasso}, \emph{logistic regression},
    \emph{random forest}\} (each method used exactly once) and justify each choice in one
    sentence:
    \begin{enumerate}[label=(\roman*),nosep,topsep=1mm]
      \item Estimate, with a $95\%$ confidence interval, how much one additional year of
        schooling raises hourly wage, from $n=3{,}000$ workers and a handful of predictors.
      \item From $n=80$ patients with $p=8{,}000$ metabolite measurements, predict a continuous
        biomarker level and identify a short list of relevant metabolites.
      \item Predict whether a customer will churn (yes/no) and report to management how each
        customer characteristic changes the \emph{odds} of churning.
      \item Forecast next-day electricity demand (continuous) from hundreds of weather and
        calendar variables with strong non-linear effects and interactions; only forecast
        accuracy matters.
    \end{enumerate}
\end{enumerate}

\ifdefined\withsolutions
\begin{solutionbox}
\textbf{(a) Bootstrap inclusion probability.}
(i) In each of the $n$ independent draws, a fixed observation is \emph{missed} with probability
$1-\tfrac1n$; since the draws are independent, all $n$ draws miss it with probability
$(1-\tfrac1n)^n$. As $n\to\infty$, $(1-\tfrac1n)^n\to e^{-1}=\mathbf{0.368}$.
(ii) The observation is therefore \emph{included} with probability $1-e^{-1}=1-0.368=0.632$, so
on average about $\mathbf{63.2\%}$ of the original observations appear in a bootstrap sample
(the remaining $\approx 1/3$ are the ``out-of-bag'' observations).
(iii) For $n=5$: $\Pr(\text{included})=1-(1-\tfrac15)^5=1-0.8^{5}=1-0.328=\mathbf{0.672}$.

\medskip
\textbf{(b) Bias--variance decomposition.}
\[
  \text{MSE}_1=\operatorname{Var}+\operatorname{Bias}^2+\operatorname{Var}(\varepsilon)
   =3.5+(-1.0)^2+2.0=3.5+1.0+2.0=\mathbf{6.5},
\]
\[
  \text{MSE}_2=0.5+(-2.5)^2+2.0=0.5+6.25+2.0=\mathbf{8.75}.
\]
Deploy the \textbf{flexible spline} (method 1): its higher estimation variance ($3.5$ vs.\
$0.5$) is more than offset by its much smaller squared bias ($1.0$ vs.\ $6.25$), giving the
lower expected test error ($6.5<8.75$). The lowest test MSE any method can achieve at $x_0$ is
the \emph{irreducible error} $\operatorname{Var}(\varepsilon)=\mathbf{2.0}$; no estimator can go
below it.

\medskip
\textbf{(c) Elastic net.}
The elastic net minimises, for $0\le\alpha\le 1$,
\[
  \sum_{i=1}^n\bigl(y_i-\beta_0-\textstyle\sum_j\beta_jx_{ij}\bigr)^2
  +\lambda\Bigl[\tfrac{1-\alpha}{2}\textstyle\sum_j\beta_j^2+\alpha\textstyle\sum_j|\beta_j|\Bigr],
\]
i.e.\ a weighted mix of the ridge ($\ell_2$) and lasso ($\ell_1$) penalties. From the
\emph{lasso} it inherits the ability to set coefficients \emph{exactly to zero} (sparsity /
variable selection); from \emph{ridge} it inherits the graceful handling of correlated
predictors and the ability to select more than $n$ variables when $p>n$. It clearly beats the
lasso when there are \emph{groups of highly correlated predictors}: the lasso tends to pick one
member of a group essentially at random and drop the rest, whereas the elastic net's ridge part
shares the weight across the group and keeps them in or out together (the ``grouping effect'').

\medskip
\textbf{(d) Collinearity and VIF.}
Collinearity does not bias the OLS estimates, but it makes them \emph{imprecise and unstable}:
it inflates their variances and standard errors, widens confidence intervals, and can make
individual coefficients change sign or lose significance even when the predictors are jointly
important (predictions may still be fine). With $R_j^2=0.9$,
\[
  \mathrm{VIF}_j=\frac{1}{1-0.9}=\frac{1}{0.1}=\mathbf{10},
\]
so the variance of $\hat\beta_j$ is $10$ times larger than it would be with an uncorrelated
$X_j$, and its standard error is inflated by $\sqrt{10}=\mathbf{3.162}$. As a rule of thumb a
$\mathrm{VIF}$ above $5$--$10$ signals problematic collinearity.

\medskip
\textbf{(e) Matching.}
\begin{itemize}[nosep,leftmargin=5mm]
  \item (i) \textbf{Linear regression}: quantitative response, goal is \emph{inference}
    (effect size with confidence interval), $n\gg p$ --- exactly what OLS with standard errors
    provides.
  \item (ii) \textbf{Lasso}: $p\gg n$ ($8{,}000$ vs.\ $80$) makes OLS infeasible; the $\ell_1$
    penalty regularises and selects a short list of metabolites.
  \item (iii) \textbf{Logistic regression}: binary response with interpretable coefficients
    ($e^{\hat\beta_j}$ = odds ratios) --- suitable for reporting how each characteristic changes
    the odds.
  \item (iv) \textbf{Random forest}: captures non-linearities and interactions automatically
    from many predictors and is accurate and robust; interpretability is not required, only
    forecast accuracy.
\end{itemize}
\end{solutionbox}
\fi

\vspace{6mm}
\begin{center}
\rule{0.35\textwidth}{0.4pt}\\[1mm]
\textit{End of the examination --- good luck!}
\end{center}

\end{document}
"
% ============================================================
\documentclass[11pt,a4paper]{article}
\usepackage[utf8]{inputenc}
\usepackage[T1]{fontenc}
\usepackage[english]{babel}
\usepackage[margin=2.2cm]{geometry}
\usepackage{amsmath,amssymb}
\usepackage{booktabs,array}
\usepackage{enumitem}
\usepackage{xcolor}
\usepackage{tcolorbox}
\tcbuselibrary{skins,breakable}

\setlength{\parindent}{0pt}

% Teal solution box (only shown when \withsolutions is defined)
\newtcolorbox{solutionbox}{enhanced,breakable,colback=teal!5,colframe=teal!55!black,
  boxrule=0pt,leftrule=2.5pt,arc=2pt,fonttitle=\bfseries,title=Solution,
  left=2mm,right=2mm,top=1mm,bottom=1.5mm}

\newcommand{\problem}[2]{\par\vspace{7mm}\noindent{\large\textbf{Problem #1}}%
  \hfill{\large\textbf{(#2 points)}}\par\nopagebreak\vspace{0.6mm}%
  \noindent\rule{\textwidth}{0.6pt}\par\nopagebreak\vspace{2.5mm}}
\newcommand{\pp}[1]{\textbf{[#1~P]}\enspace}

\begin{document}

% ------------------------------------------------------------
% Header block
% ------------------------------------------------------------
\begin{center}
  {\Large\textbf{HSBI --- Bielefeld University of Applied Sciences and Arts}}\\[1.2mm]
  {\large Quantitative Research Methods}\\[1mm]
  {\normalsize Summer Semester 2026 \quad$\cdot$\quad Examiner: Prof.\ Dr.\ Christoph Weisser}\\[4.5mm]
  {\LARGE\textbf{Final Mock Examination --- Set C (Lectures 1--12)}}\\[1.2mm]
  {\normalsize Comprehensive coverage of ISLP Chapters 2--8, 10 and 13,\\
   with emphasis on Chapters 7 (splines/GAMs), 8 (trees), 10 (deep learning) and 13 (multiple testing)}%
  \ifdefined\withsolutions\\[2.2mm]{\large\textcolor{teal!55!black}{\textbf{--- Version with detailed solutions ---}}}\fi
\end{center}

\vspace{2.5mm}
\begin{center}
\begin{tabular}{@{}ll@{}}
  \textbf{Duration:} & \textbf{120 minutes}\\
  \textbf{Total credit:} & \textbf{120 points}\\
  \textbf{Allowed aids:} & non-programmable calculator; one handwritten A4 sheet of notes (both sides)\\
\end{tabular}
\end{center}

\vspace{3mm}
Name: \rule[-1mm]{0.38\textwidth}{0.4pt}\hfill Matriculation no.: \rule[-1mm]{0.24\textwidth}{0.4pt}

\vspace{4.5mm}
\begin{center}
\renewcommand{\arraystretch}{1.35}
\begin{tabular}{|l|c|c|c|c|c|c|c|c|}
\hline
\textbf{Problem} & \textbf{P1} & \textbf{P2} & \textbf{P3} & \textbf{P4} & \textbf{P5} & \textbf{P6} & \textbf{P7} & \textbf{Total}\\\hline
Maximum points & 16 & 8 & 20 & 16 & 20 & 20 & 20 & \textbf{120}\\\hline
Achieved points & \hspace{8mm} & \hspace{8mm} & \hspace{8mm} & \hspace{8mm} & \hspace{8mm} & \hspace{8mm} & \hspace{8mm} & \hspace{10mm}\\\hline
\end{tabular}
\end{center}

\vspace{2.5mm}
\textbf{Instructions}
\begin{itemize}[nosep,leftmargin=5.5mm]
  \item Justify all answers. Numerical results without a derivation or a brief reasoning receive no credit.
  \item Round final numerical results to 3 decimal places unless stated otherwise. Small deviations caused by carrying rounded intermediate results are accepted.
  \item All numerical constants you need (powers of $e$, logarithms, $(0.95)^{40}$, $t$-quantiles) are provided in the problems.
  \item The number of points per problem roughly equals the recommended working time in minutes.
\end{itemize}
\vspace{1mm}\noindent\rule{\textwidth}{0.6pt}

% ============================================================
\problem{1: Moving beyond linearity --- splines}{16}
{\small\itshape Lecture slides: Chapter 7 --- Moving Beyond Linearity (Lecture 9).}\par\vspace{1mm}

\begin{enumerate}[label=(\alph*),leftmargin=7mm,itemsep=2.5mm]
  \item \pp{4} A \emph{cubic spline} is a piecewise cubic polynomial that is continuous and has
    continuous first and second derivatives at each knot. Derive the number of degrees of freedom
    (free parameters, counting the intercept) of a cubic spline with $K$ interior knots by
    counting the coefficients of the polynomial pieces and subtracting the constraints imposed
    at the knots. Evaluate your formula for $K=3$ interior knots.
  \item \pp{4} How many degrees of freedom does a \emph{natural} cubic spline with the same
    $K=3$ knots use? State the additional constraint a natural spline imposes, say in which
    region of the predictor space it takes effect, and explain why it is useful in practice.
  \item \pp{4} A \emph{smoothing spline} is the function $g$ that minimises
    \[
      \sum_{i=1}^{n}\bigl(y_i-g(x_i)\bigr)^2+\lambda\int g''(t)^2\,dt .
    \]
    Explain the role of each of the two terms, and describe the behaviour of the solution
    $\hat{g}$ in the two limiting cases $\lambda\to 0$ and $\lambda\to\infty$.
  \item \pp{4} Write down the \emph{truncated-power basis} representation of a cubic spline with
    a single knot at $\xi$, define the truncated-power basis function explicitly, and explain in
    one sentence why adding this basis function does not destroy the smoothness of the fit at
    $\xi$.
\end{enumerate}

\ifdefined\withsolutions
\begin{solutionbox}
\textbf{(a) Degrees of freedom of a cubic spline.}
$K$ interior knots split the range of $X$ into $K+1$ intervals, each carrying a cubic
polynomial with $4$ coefficients: $4(K+1)=4K+4$ parameters. At each knot, continuity of
$g$, $g'$ and $g''$ imposes $3$ constraints, i.e.\ $3K$ constraints in total. Hence
\[
  \mathrm{df}=4(K+1)-3K=K+4 .
\]
For $K=3$: $\mathrm{df}=3+4=\mathbf{7}$ free parameters (equivalently: intercept, $x$, $x^2$,
$x^3$ plus one truncated-power basis function per knot $=4+K$).

\medskip
\textbf{(b) Natural cubic spline.}
A natural spline additionally requires the function to be \emph{linear beyond the boundary
knots}, i.e.\ in the two outer regions $X<\xi_1$ and $X>\xi_K$. Forcing the quadratic and cubic
terms to vanish there removes $2$ parameters per boundary region, i.e.\ $4$ in total:
\[
  \mathrm{df}=(K+4)-4=K,\qquad K=3:\ \mathrm{df}=\mathbf{3}.
\]
Usefulness: polynomials and unconstrained splines are notoriously erratic near the boundaries,
where data are sparse; the linearity constraint stabilises the fit there
(\emph{lower variance} of predictions at and beyond the boundary), at the price of a small
amount of additional bias.

\medskip
\textbf{(c) Smoothing spline.}
The first term (RSS) rewards fidelity to the data; the second term penalises roughness, since
$g''$ measures the curvature (wiggliness) of $g$, and $\lambda\ge 0$ is the tuning parameter
that trades the two off.
\emph{$\lambda\to 0$:} the penalty disappears; $\hat{g}$ can be made arbitrarily wiggly and
\emph{interpolates} the training observations exactly (zero training RSS, severe overfitting;
effective df $\to n$).
\emph{$\lambda\to\infty$:} any curvature is infinitely expensive, so $g''\equiv 0$ is enforced;
$\hat{g}$ becomes a straight line and equals the \emph{least-squares regression line}
(effective df $\to 2$). Between the extremes, $\lambda$ (chosen e.g.\ by LOOCV) controls the
effective degrees of freedom, which decrease smoothly from $n$ to $2$.

\medskip
\textbf{(d) Truncated-power basis with one knot $\xi$.}
\[
  f(x)=\beta_0+\beta_1x+\beta_2x^2+\beta_3x^3+\beta_4\,(x-\xi)^3_+,
  \qquad
  (x-\xi)^3_+=\begin{cases}(x-\xi)^3, & x>\xi\\[0.5mm] 0, & x\le\xi\end{cases}
\]
The function $(x-\xi)^3_+$ and its first and second derivatives are all $0$ at $x=\xi$, so
adding it changes only the third derivative at the knot; the fit remains continuous with
continuous first and second derivatives --- exactly the defining smoothness of a cubic spline.
\end{solutionbox}
\fi

% ============================================================
\problem{2: Generalised additive models}{8}
{\small\itshape Lecture slides: Chapter 7 --- Moving Beyond Linearity (Lecture 9).}\par\vspace{1mm}

\begin{enumerate}[label=(\alph*),leftmargin=7mm,itemsep=2.5mm]
  \item \pp{3} Write down the general form of a \emph{generalised additive model} (GAM) for a
    quantitative response $Y$ and predictors $X_1,\dots,X_p$, and explain in one sentence how it
    generalises multiple linear regression.
  \item \pp{3} State one important \emph{advantage} of GAMs concerning the interpretation of the
    fitted functions $f_j$, and the main structural \emph{limitation} of the additive form ---
    including how this limitation can be repaired.
  \item \pp{2} Describe in one or two sentences how the \emph{backfitting} algorithm fits a GAM.
\end{enumerate}

\ifdefined\withsolutions
\begin{solutionbox}
\textbf{(a) Model form.}
\[
  Y=\beta_0+f_1(X_1)+f_2(X_2)+\dots+f_p(X_p)+\varepsilon .
\]
Multiple linear regression is the special case $f_j(X_j)=\beta_jX_j$; a GAM replaces each linear
component by a smooth non-linear function $f_j$ (e.g.\ a natural or smoothing spline, or local
regression), while keeping the \emph{additive} structure.

\medskip
\textbf{(b) Advantage and limitation.}
\emph{Advantage (interpretability):} because the model is additive, the contribution of each
$X_j$ to $Y$ is captured by its own function $f_j$, which can be examined and plotted
individually while the other variables are held fixed --- non-linear effects remain as easy to
read as in a coefficient table, one panel per predictor.
\emph{Limitation:} additivity excludes \emph{interactions}: the effect of $X_j$ cannot depend
on the value of $X_k$. Repair: add interaction terms manually, e.g.\ a product term
$X_j\times X_k$ or a low-dimensional smooth surface $f_{jk}(X_j,X_k)$ as an extra component.

\medskip
\textbf{(c) Backfitting.}
Cycle through the predictors and repeatedly re-fit each $f_j$ (by a smoother, e.g.\ a spline)
to the \emph{partial residuals} $r_i=y_i-\beta_0-\sum_{k\ne j}f_k(x_{ik})$, holding all other
fitted functions fixed; iterate until the $f_j$ stop changing.
\end{solutionbox}
\fi

% ============================================================
\problem{3: Regression and classification trees by hand}{20}
{\small\itshape Lecture slides: Chapter 8 --- Tree-Based Methods (Lecture 10).}\par\vspace{1mm}

A logistics company records for $n=6$ regional depots the number of delivery vans in operation
$x$ and the number of parcels delivered that day $y$ (in 1{,}000 parcels):

\begin{center}
\begin{tabular}{lcccccc}
\toprule
Depot $i$ & 1 & 2 & 3 & 4 & 5 & 6\\
\midrule
Vans $x_i$ & 1 & 2 & 3 & 4 & 5 & 6\\
Parcels $y_i$ & 3 & 5 & 4 & 9 & 11 & 10\\
\bottomrule
\end{tabular}
\end{center}

A regression tree with a \emph{single} split at cutpoint $s$ partitions the data into
$R_1=\{i: x_i<s\}$ and $R_2=\{i: x_i\ge s\}$ and predicts the region mean $\hat{y}_{R_m}$ in
each region.

\begin{enumerate}[label=(\alph*),leftmargin=7mm,itemsep=2.5mm]
  \item \pp{8} For each of the two candidate cutpoints $s=2.5$ and $s=3.5$, determine the two
    regions, the region means and the resulting total residual sum of squares
    $\operatorname{RSS}(s)=\sum_{m=1}^{2}\sum_{i\in R_m}(y_i-\hat{y}_{R_m})^2$.
    Which cutpoint does recursive binary splitting choose, and why?
  \item \pp{2} Using the chosen split, predict the number of parcels for a new depot that will
    operate $x=5$ vans.
  \item \pp{6} A \emph{classification} tree for a different task has a terminal node containing
    $9$ training observations from three product categories: $5$ ``Electronics'', $3$
    ``Clothing'' and $1$ ``Books''. Compute the class proportions, the Gini index
    $G=\sum_{k}\hat{p}_{k}(1-\hat{p}_{k})$ and the entropy $D=-\sum_{k}\hat{p}_{k}\ln\hat{p}_{k}$
    of this node \emph{(use $\ln(5/9)=-0.5878$, $\ln(1/3)=-1.0986$ and $\ln(1/9)=-2.1972$)}.
    Which class does the node predict?
  \item \pp{4} In \emph{cost-complexity pruning} one minimises, for a subtree
    $T\subseteq T_0$ of the large tree $T_0$,
    \[
      \sum_{m=1}^{|T|}\ \sum_{i:\,x_i\in R_m}\bigl(y_i-\hat{y}_{R_m}\bigr)^2+\alpha\,|T| .
    \]
    Explain the role of the tuning parameter $\alpha$: what happens for $\alpha=0$, what happens
    as $\alpha$ increases, and how is $\alpha$ chosen in practice?
\end{enumerate}

\ifdefined\withsolutions
\begin{solutionbox}
\textbf{(a) Evaluating the two cutpoints.}
Root node for reference: $\bar{y}=42/6=7$.

\emph{Cutpoint $s=2.5$:} $R_1=\{1,2\}$ with $y$-values $\{3,5\}$, mean $\hat{y}_{R_1}=4$;
$R_2=\{3,4,5,6\}$ with $y$-values $\{4,9,11,10\}$, mean $\hat{y}_{R_2}=34/4=8.5$.
\begin{align*}
  \operatorname{RSS}_{R_1}&=(3-4)^2+(5-4)^2=2,\\
  \operatorname{RSS}_{R_2}&=(4-8.5)^2+(9-8.5)^2+(11-8.5)^2+(10-8.5)^2=20.25+0.25+6.25+2.25=29,
\end{align*}
so $\operatorname{RSS}(2.5)=2+29=\mathbf{31}$.

\emph{Cutpoint $s=3.5$:} $R_1=\{1,2,3\}$ with $y$-values $\{3,5,4\}$, mean
$\hat{y}_{R_1}=12/3=4$; $R_2=\{4,5,6\}$ with $y$-values $\{9,11,10\}$, mean
$\hat{y}_{R_2}=30/3=10$.
\[
  \operatorname{RSS}_{R_1}=(3-4)^2+(5-4)^2+(4-4)^2=2,\qquad
  \operatorname{RSS}_{R_2}=(9-10)^2+(11-10)^2+(10-10)^2=2,
\]
so $\operatorname{RSS}(3.5)=2+2=\mathbf{4}$.

Since $4<31$, recursive binary splitting is \emph{greedy} and picks the cutpoint with the
smallest RSS: $\mathbf{s=3.5}$. (It cleanly separates the three low-volume depots, mean $4$,
from the three high-volume depots, mean $10$; in fact $s=3.5$ is the best of all five possible
cutpoints.)

\medskip
\textbf{(b) Prediction for $x=5$.}
$x=5\ge 3.5$, so the new depot falls into $R_2$ and the tree predicts
$\hat{y}=\hat{y}_{R_2}=\mathbf{10}$ (i.e.\ $10{,}000$ parcels).

\medskip
\textbf{(c) Node impurity (three classes).}
$\hat{p}_{\text{E}}=5/9$, $\hat{p}_{\text{C}}=3/9=1/3$, $\hat{p}_{\text{B}}=1/9$.
\[
  G=\tfrac{5}{9}\cdot\tfrac{4}{9}+\tfrac{3}{9}\cdot\tfrac{6}{9}+\tfrac{1}{9}\cdot\tfrac{8}{9}
   =\tfrac{20+18+8}{81}=\tfrac{46}{81}=\mathbf{0.568}.
\]
\[
  D=-\Bigl(\tfrac{5}{9}\ln\tfrac{5}{9}+\tfrac{1}{3}\ln\tfrac{1}{3}+\tfrac{1}{9}\ln\tfrac{1}{9}\Bigr)
   =-\bigl(\tfrac{5}{9}(-0.5878)+\tfrac{1}{3}(-1.0986)+\tfrac{1}{9}(-2.1972)\bigr)
\]
\[
   =0.3266+0.3662+0.2441=\mathbf{0.937}.
\]
The node predicts the majority class, \textbf{Electronics} (with estimated probability
$5/9\approx 0.556$). Both $G$ and $D$ would be $0$ for a pure node; here the node is quite
impure because no class dominates strongly.

\medskip
\textbf{(d) Role of $\alpha$.}
$\alpha$ is the price per terminal node: it trades off the fit of the subtree (first term)
against its complexity $|T|$.
For $\alpha=0$ the criterion equals the training RSS and the full tree $T_0$ is optimal.
As $\alpha$ increases, additional leaves must ``pay for themselves''; branches whose RSS
reduction is smaller than $\alpha$ are pruned, so the optimal subtree becomes smaller ---
one obtains a \emph{nested sequence} of subtrees as $\alpha$ grows (weakest-link pruning).
In practice $\alpha$ is selected by \emph{cross-validation} (choose the $\alpha$ with the
smallest CV error, then refit on all data), which balances bias against variance instead of
optimising training fit.
\end{solutionbox}
\fi

% ============================================================
\problem{4: Ensembles --- bagging, random forests, boosting}{16}
{\small\itshape Lecture slides: Chapter 8 --- Tree-Based Methods (Lecture 10).}\par\vspace{1mm}

\begin{enumerate}[label=(\alph*),leftmargin=7mm,itemsep=2.5mm]
  \item \pp{6} Let $\hat{f}_1(x),\dots,\hat{f}_B(x)$ be $B$ identically distributed (but not
    independent) tree predictions at a fixed point $x$, each with variance $\sigma^2$ and with
    pairwise correlation $\rho$. The variance of the bagged average is
    \[
      \operatorname{Var}\Bigl(\tfrac{1}{B}\sum_{b=1}^{B}\hat{f}_b(x)\Bigr)
      =\rho\,\sigma^{2}+\frac{1-\rho}{B}\,\sigma^{2}.
    \]
    (i) Explain briefly why this formula shows that bagging reduces variance, and which term
    limits the reduction. (ii) Evaluate the variance for $B=50$, $\rho=0.4$ and $\sigma^2=16$,
    and compare it with the variance of a single tree. (iii) What value does the variance
    approach as $B\to\infty$?
  \item \pp{4} How do \emph{random forests} modify bagging at each split, and what is the
    recommended subset size $m$ for \emph{regression} as a function of $p$ (evaluate for
    $p=30$)? Explain, with reference to the formula in (a), why this modification improves on
    bagging.
  \item \pp{4} Name the three tuning parameters of \emph{boosting} for regression trees and
    state in one sentence each what they control. What is the ``slow learning'' idea behind
    boosting?
  \item \pp{2} What does the \emph{out-of-bag} (OOB) error of a bagged model estimate, and why
    is it available essentially for free?
\end{enumerate}

\ifdefined\withsolutions
\begin{solutionbox}
\textbf{(a) Variance of the bagged average.}
(i) The second term $\frac{1-\rho}{B}\sigma^2$ is driven to $0$ by averaging over many trees
($B$ large), so averaging removes the ``independent part'' of the trees' variability ---
this is the variance reduction of bagging. The first term $\rho\sigma^2$ does \emph{not}
depend on $B$: because the bootstrap trees are grown on overlapping samples of the same data
(and share the same strong predictors), they are positively correlated, and this correlated
part of the variance cannot be averaged away. It limits the reduction.
(ii) For $B=50$, $\rho=0.4$, $\sigma^2=16$:
\[
  \operatorname{Var}=0.4\cdot 16+\frac{1-0.4}{50}\cdot 16
  =6.4+\frac{0.6\cdot 16}{50}=6.4+0.192=\mathbf{6.592},
\]
compared with $\sigma^2=16$ for a single tree --- a reduction of about $58.8\%$.
(iii) As $B\to\infty$ the variance approaches the floor $\rho\sigma^2=0.4\cdot 16=\mathbf{6.4}$.
(Increasing $B$ further does not overfit; it only stabilises the average.)

\medskip
\textbf{(b) Random forests.}
At each split, only a \emph{random subset} of $m$ of the $p$ predictors is considered as
splitting candidates; a fresh subset is drawn at every split. Recommended for regression:
$m\approx p/3$, so for $p=30$: $m=30/3=\mathbf{10}$.
Why it helps: with bagging, a dominant predictor is used in the top split of almost every tree,
so the trees look alike and $\rho$ is high. Excluding it from $\approx (p-m)/p$ of the split
searches \emph{decorrelates} the trees, i.e.\ lowers $\rho$ --- and by (a) this lowers the
irreducible floor $\rho\sigma^2$ of the ensemble variance, which mere averaging cannot touch.

\medskip
\textbf{(c) Boosting.}
\begin{itemize}[nosep,leftmargin=5mm]
  \item \emph{Number of trees $B$:} how many small trees are added sequentially; unlike
    bagging, boosting \emph{can} overfit if $B$ is too large, so $B$ is chosen by
    cross-validation.
  \item \emph{Shrinkage $\lambda$} (learning rate, e.g.\ $0.01$ or $0.001$): scales the
    contribution of each new tree and thus the speed of learning; smaller $\lambda$ usually
    needs larger $B$.
  \item \emph{Number of splits $d$ per tree} (interaction depth): the complexity of each tree
    and the interaction order of the model; often even $d=1$ (stumps, additive model) works
    well.
\end{itemize}
\emph{Slow learning:} each tree is fitted to the current \emph{residuals}, and only a small
fraction $\lambda$ of it is added to the model; the model thus improves slowly and exactly in
the regions where it still performs poorly, instead of fitting the data hard in one step ---
which tends to generalise better.

\medskip
\textbf{(d) OOB error.}
Each bootstrap sample leaves out on average about $1/3$ of the observations
(``out-of-bag''). Predicting each observation only with the trees that did \emph{not} use it,
and aggregating these predictions, yields the OOB error --- a valid estimate of the
\emph{test error} of the bagged model. It is free because it reuses the bootstrap side
product; no separate validation set or cross-validation loop is required.
\end{solutionbox}
\fi

% ============================================================
\problem{5: Neural networks by hand}{20}
{\small\itshape Lecture slides: Chapter 10 --- Deep Learning (Lecture 11).}\par\vspace{1mm}

Consider a tiny feed-forward network for a quantitative response with input
$x=(x_1,x_2)=(2,3)$, one hidden layer with two ReLU units ($g(z)=\max\{0,z\}$) and a linear
output unit:
\[
  A_1=g\bigl(-3+2\,x_1+1\,x_2\bigr),\qquad
  A_2=g\bigl(2-2\,x_1-1\,x_2\bigr),\qquad
  f(x)=-2+3\,A_1+2\,A_2 .
\]

\begin{enumerate}[label=(\alph*),leftmargin=7mm,itemsep=2.5mm]
  \item \pp{8} Perform the forward pass step by step: compute both pre-activations, both
    activations $A_1,A_2$ and the network output $f(x)$. Comment briefly on what the ReLU did
    to the second unit and why such non-linear activations are essential.
  \item \pp{6} A classification network with $K=3$ classes produces the logits (output-layer
    pre-activations) $z=(z_1,z_2,z_3)=(2.0,\ 0.0,\ -1.0)$ for one observation whose true class
    is class~2. Compute all three softmax probabilities and the cross-entropy loss
    $-\ln \hat{p}_{\text{true}}$.
    \emph{(Use $e^{2}=7.389$, $e^{0}=1.000$, $e^{-1}=0.368$ and $\ln 8.757=2.170$.)}
  \item \pp{4} Count the parameters of a fully connected network with $p=6$ inputs, one hidden
    layer with $k=5$ units and an output layer with $K=3$ units, where every hidden and output
    unit has a bias. Show the count separately for the two weight layers.
  \item \pp{2} A convolutional layer receives a $64\times 64$ image and uses filters of size
    $F=3$ with padding $P=1$ and stride $S=2$. Compute the spatial size of the output feature
    map from $\bigl\lfloor(W-F+2P)/S\bigr\rfloor+1$ (note the stride is $2$, so the division is
    rounded down).
\end{enumerate}

\ifdefined\withsolutions
\begin{solutionbox}
\textbf{(a) Forward pass.}
Pre-activations:
\[
  z_1=-3+2\cdot 2+1\cdot 3=-3+4+3=4,\qquad
  z_2=2-2\cdot 2-1\cdot 3=2-4-3=-5 .
\]
ReLU activations:
\[
  A_1=g(4)=\max\{0,4\}=4,\qquad A_2=g(-5)=\max\{0,-5\}=0 .
\]
Output:
\[
  f(x)=-2+3\cdot 4+2\cdot 0=-2+12+0=\mathbf{10}.
\]
The ReLU \emph{switched the second unit off}: its negative pre-activation is clipped to $0$,
so unit 2 contributes nothing for this input (but would for other inputs). This thresholding
is the source of non-linearity: without a non-linear $g$, the network would collapse into a
single linear function of $x$, no matter how many layers it has.

\medskip
\textbf{(b) Softmax and cross-entropy.}
Exponentials: $e^{2.0}=7.389$, $e^{0.0}=1.000$, $e^{-1.0}=0.368$; their sum is
$7.389+1.000+0.368=8.757$. Softmax probabilities:
\[
  \hat{p}_1=\frac{7.389}{8.757}=0.844,\qquad
  \hat{p}_2=\frac{1.000}{8.757}=0.114,\qquad
  \hat{p}_3=\frac{0.368}{8.757}=0.042,
\]
which sum to $1.000$. Cross-entropy loss for the true class $2$:
\[
  -\ln\hat{p}_2=-\ln\frac{1.000}{8.757}=-\bigl(0-\ln 8.757\bigr)=\ln 8.757=\mathbf{2.170}.
\]
(Note that the model's most probable class is class~1, but the loss is evaluated at the
\emph{true} class~2, which received only probability $0.114$ --- hence the sizeable loss.)

\medskip
\textbf{(c) Parameter count for $6\to 5\to 3$.}
Input $\to$ hidden: $6\cdot 5=30$ weights $+\,5$ biases $=35$.
Hidden $\to$ output: $5\cdot 3=15$ weights $+\,3$ biases $=18$.
Total: $35+18=\mathbf{53}$ parameters.

\medskip
\textbf{(d) Convolution output size.}
\[
  \Bigl\lfloor\frac{W-F+2P}{S}\Bigr\rfloor+1
  =\Bigl\lfloor\frac{64-3+2\cdot 1}{2}\Bigr\rfloor+1
  =\Bigl\lfloor\frac{63}{2}\Bigr\rfloor+1=\lfloor 31.5\rfloor+1=31+1=32,
\]
so the output feature map is $\mathbf{32\times 32}$ --- the stride $S=2$ roughly halves the
$64\times 64$ input (the floor discards the fractional part before adding $1$).
\end{solutionbox}
\fi

% ============================================================
\problem{6: Multiple testing}{20}
{\small\itshape Lecture slides: Chapter 13 --- Multiple Testing (Lecture 12).}\par\vspace{1mm}

\begin{enumerate}[label=(\alph*),leftmargin=7mm,itemsep=2.5mm]
  \item \pp{4} A lab performs $m=40$ \emph{independent} hypothesis tests, each at level
    $\alpha=0.05$, and suppose all $40$ null hypotheses are true. Define the family-wise error
    rate (FWER) and compute it for this situation \emph{(use $(0.95)^{40}=0.1285$)}. Interpret
    the result in one sentence.
  \item \pp{7} A research group screens $m=6$ candidate biomarkers; the ordered $p$-values are
    \begin{center}
    \begin{tabular}{lcccccc}
    \toprule
    $j$ & 1 & 2 & 3 & 4 & 5 & 6\\
    \midrule
    $p_{(j)}$ & 0.001 & 0.008 & 0.011 & 0.030 & 0.041 & 0.260\\
    \bottomrule
    \end{tabular}
    \end{center}
    Control the FWER at $\alpha=0.05$ (i) with \emph{Bonferroni} and (ii) with the
    \emph{Holm step-down} procedure. Give the relevant thresholds, state for each method
    exactly which hypotheses are rejected, and explain why Holm can never reject fewer
    hypotheses than Bonferroni.
  \item \pp{6} Apply the \emph{Benjamini--Hochberg} (BH) procedure with FDR level $q=0.05$ to
    the same six $p$-values: give the thresholds $jq/m$, determine which hypotheses are
    rejected, and state precisely what ``FDR control at $q=0.05$'' guarantees for the
    resulting set of rejections.
  \item \pp{3} A pharmaceutical company runs a \emph{confirmatory} phase-III trial that tests
    $m=5$ pre-specified endpoints; a single false-positive endpoint could lead to the approval
    of an ineffective drug. Should they control the FWER or the FDR? Justify your
    recommendation.
\end{enumerate}

\ifdefined\withsolutions
\begin{solutionbox}
\textbf{(a) FWER for $m=40$ independent tests.}
$\text{FWER}=\Pr(\text{at least one Type I error})=\Pr(V\ge 1)$, where $V$ is the number of
true nulls that are (falsely) rejected. With independence and all nulls true:
\[
  \text{FWER}=1-\Pr(\text{no false rejection})=1-(1-0.05)^{40}=1-(0.95)^{40}
  =1-0.1285=\mathbf{0.8715}.
\]
Interpretation: even though each single test is performed at the $5\%$ level, there is about an
$87\%$ chance of at least one false discovery among the $40$ tests --- testing many hypotheses
without adjustment almost guarantees spurious findings.

\medskip
\textbf{(b) Bonferroni and Holm at $\alpha=0.05$.}
\emph{Bonferroni:} reject $H_{(j)}$ iff $p_{(j)}\le\alpha/m=0.05/6=0.00833$.
$p_{(1)}=0.001\le 0.00833$ \checkmark, $p_{(2)}=0.008\le 0.00833$ \checkmark,
$p_{(3)}=0.011>0.00833$ $\times$.
$\Rightarrow$ Bonferroni rejects $H_{(1)}$ and $H_{(2)}$: \textbf{2 rejections}.

\emph{Holm:} compare $p_{(j)}$ with $\alpha/(m-j+1)$ step by step and stop at the first failure:
\begin{center}
\begin{tabular}{lcccccc}
\toprule
$j$ & 1 & 2 & 3 & 4 & 5 & 6\\
\midrule
$p_{(j)}$ & 0.001 & 0.008 & 0.011 & 0.030 & 0.041 & 0.260\\
$\alpha/(m-j+1)$ & 0.00833 & 0.0100 & 0.0125 & 0.0167 & 0.0250 & 0.0500\\
reject? & \checkmark & \checkmark & \checkmark & stop & --- & ---\\
\bottomrule
\end{tabular}
\end{center}
$0.001\le 0.00833$, $0.008\le 0.0100$, $0.011\le 0.0125$, but $0.030>0.0167$: stop.
$\Rightarrow$ Holm rejects $H_{(1)},H_{(2)},H_{(3)}$: \textbf{3 rejections}.

\emph{Why Holm $\supseteq$ Bonferroni:} Holm's thresholds $\alpha/(m-j+1)$ are $\ge\alpha/m$
for every $j$ (with equality only at $j=1$), so every $p$-value that passes Bonferroni also
passes Holm's less stringent later hurdles --- Holm is uniformly more powerful, yet still
controls the FWER at $\alpha$ without any independence assumption. Here Holm rejects one more
hypothesis ($H_{(3)}$) than Bonferroni.

\medskip
\textbf{(c) Benjamini--Hochberg at $q=0.05$.}
Thresholds $jq/m=j\cdot 0.05/6$:
\begin{center}
\begin{tabular}{lcccccc}
\toprule
$j$ & 1 & 2 & 3 & 4 & 5 & 6\\
\midrule
$p_{(j)}$ & 0.001 & 0.008 & 0.011 & 0.030 & 0.041 & 0.260\\
$jq/m$ & 0.00833 & 0.0167 & 0.0250 & 0.0333 & 0.0417 & 0.0500\\
$p_{(j)}\le jq/m$? & yes & yes & yes & yes & \textbf{yes} & no\\
\bottomrule
\end{tabular}
\end{center}
$L=$ largest $j$ with $p_{(j)}\le jq/m$: $j=5$ (since $0.041\le 0.0417$, while
$0.260>0.0500$). BH rejects \emph{all} $H_{(j)}$ with $j\le L$: $H_{(1)},\dots,H_{(5)}$:
\textbf{5 rejections}.
Guarantee: $\text{FDR}=E\bigl[V/\max(R,1)\bigr]\le q=0.05$ --- \emph{on average} at most
$5\%$ of the rejected hypotheses are false discoveries. It does \emph{not} guarantee that at
most $5\%$ of these particular $5$ rejections are false, nor does it bound the probability of
making any false rejection (that would be FWER control). Note BH is the most liberal of the
three here ($5$ rejections vs.\ $3$ for Holm and $2$ for Bonferroni).

\medskip
\textbf{(d) Confirmatory phase-III trial.}
Control the \textbf{FWER}. This is a confirmatory study with only $m=5$ pre-specified
endpoints, and the cost of even a single false positive is severe (approval of an ineffective
drug, patient harm, regulatory failure). FWER control (e.g.\ Holm at $\alpha=0.05$) bounds the
probability of \emph{any} false rejection at $5\%$, which is exactly the guarantee a regulator
demands. FDR control would only limit the \emph{expected proportion} of false endpoints and
would tolerate an occasional false positive --- acceptable in large exploratory screens, but
not when each individual claim must stand on its own.
\end{solutionbox}
\fi

% ============================================================
\problem{7: Cumulative essentials (Chapters 2--6)}{20}
{\small\itshape Lecture slides: Chapters 2--6 (Lectures 1--8).}\par\vspace{1mm}

\begin{enumerate}[label=(\alph*),leftmargin=7mm,itemsep=2.5mm]
  \item \pp{4} The \emph{bootstrap} draws a resample of size $n$ with replacement from the $n$
    observations. (i) Show that the probability that a given observation is \emph{not} selected
    into a particular bootstrap sample equals $(1-1/n)^n$, and give its limit as $n\to\infty$
    \emph{(use $e^{-1}=0.368$)}. (ii) Hence what fraction of the original observations appears
    (at least once) in a typical bootstrap sample as $n\to\infty$? (iii) Compute this inclusion
    probability exactly for $n=5$ \emph{(use $0.8^{5}=0.328$)}.
  \item \pp{4} At a fixed test point $x_0$ the expected test error of an estimate $\hat f$
    decomposes as
    $E\bigl[(y_0-\hat f(x_0))^2\bigr]
      =\operatorname{Var}(\hat f(x_0))+\bigl[\operatorname{Bias}(\hat f(x_0))\bigr]^2
       +\operatorname{Var}(\varepsilon)$.
    Two candidate methods give, at $x_0$:
    \begin{center}
    \begin{tabular}{lccc}
    \toprule
    Method & $\operatorname{Var}(\hat f)$ & $\operatorname{Bias}(\hat f)$ & $\operatorname{Var}(\varepsilon)$\\
    \midrule
    (1) flexible spline & 3.5 & $-1.0$ & 2.0\\
    (2) linear model & 0.5 & $-2.5$ & 2.0\\
    \bottomrule
    \end{tabular}
    \end{center}
    Compute the expected test MSE of each method, say which you would deploy at $x_0$, and state
    the lowest test MSE \emph{any} method could achieve here.
  \item \pp{3} Write down the \emph{elastic net} penalty and explain what it inherits from the
    lasso and what it inherits from ridge regression. In which situation does the elastic net
    clearly outperform the lasso?
  \item \pp{3} Two predictors in a regression have correlation $0.95$. Explain what such
    \emph{collinearity} does to the least-squares coefficient estimates. The variance inflation
    factor is $\mathrm{VIF}_j=1/(1-R_j^2)$, where $R_j^2$ is from regressing $X_j$ on the other
    predictors; compute and interpret $\mathrm{VIF}_j$ for $R_j^2=0.9$ \emph{(use
    $\sqrt{10}=3.162$)}.
  \item \pp{6} Match each scenario to exactly one method from
    \{\emph{linear regression}, \emph{lasso}, \emph{logistic regression},
    \emph{random forest}\} (each method used exactly once) and justify each choice in one
    sentence:
    \begin{enumerate}[label=(\roman*),nosep,topsep=1mm]
      \item Estimate, with a $95\%$ confidence interval, how much one additional year of
        schooling raises hourly wage, from $n=3{,}000$ workers and a handful of predictors.
      \item From $n=80$ patients with $p=8{,}000$ metabolite measurements, predict a continuous
        biomarker level and identify a short list of relevant metabolites.
      \item Predict whether a customer will churn (yes/no) and report to management how each
        customer characteristic changes the \emph{odds} of churning.
      \item Forecast next-day electricity demand (continuous) from hundreds of weather and
        calendar variables with strong non-linear effects and interactions; only forecast
        accuracy matters.
    \end{enumerate}
\end{enumerate}

\ifdefined\withsolutions
\begin{solutionbox}
\textbf{(a) Bootstrap inclusion probability.}
(i) In each of the $n$ independent draws, a fixed observation is \emph{missed} with probability
$1-\tfrac1n$; since the draws are independent, all $n$ draws miss it with probability
$(1-\tfrac1n)^n$. As $n\to\infty$, $(1-\tfrac1n)^n\to e^{-1}=\mathbf{0.368}$.
(ii) The observation is therefore \emph{included} with probability $1-e^{-1}=1-0.368=0.632$, so
on average about $\mathbf{63.2\%}$ of the original observations appear in a bootstrap sample
(the remaining $\approx 1/3$ are the ``out-of-bag'' observations).
(iii) For $n=5$: $\Pr(\text{included})=1-(1-\tfrac15)^5=1-0.8^{5}=1-0.328=\mathbf{0.672}$.

\medskip
\textbf{(b) Bias--variance decomposition.}
\[
  \text{MSE}_1=\operatorname{Var}+\operatorname{Bias}^2+\operatorname{Var}(\varepsilon)
   =3.5+(-1.0)^2+2.0=3.5+1.0+2.0=\mathbf{6.5},
\]
\[
  \text{MSE}_2=0.5+(-2.5)^2+2.0=0.5+6.25+2.0=\mathbf{8.75}.
\]
Deploy the \textbf{flexible spline} (method 1): its higher estimation variance ($3.5$ vs.\
$0.5$) is more than offset by its much smaller squared bias ($1.0$ vs.\ $6.25$), giving the
lower expected test error ($6.5<8.75$). The lowest test MSE any method can achieve at $x_0$ is
the \emph{irreducible error} $\operatorname{Var}(\varepsilon)=\mathbf{2.0}$; no estimator can go
below it.

\medskip
\textbf{(c) Elastic net.}
The elastic net minimises, for $0\le\alpha\le 1$,
\[
  \sum_{i=1}^n\bigl(y_i-\beta_0-\textstyle\sum_j\beta_jx_{ij}\bigr)^2
  +\lambda\Bigl[\tfrac{1-\alpha}{2}\textstyle\sum_j\beta_j^2+\alpha\textstyle\sum_j|\beta_j|\Bigr],
\]
i.e.\ a weighted mix of the ridge ($\ell_2$) and lasso ($\ell_1$) penalties. From the
\emph{lasso} it inherits the ability to set coefficients \emph{exactly to zero} (sparsity /
variable selection); from \emph{ridge} it inherits the graceful handling of correlated
predictors and the ability to select more than $n$ variables when $p>n$. It clearly beats the
lasso when there are \emph{groups of highly correlated predictors}: the lasso tends to pick one
member of a group essentially at random and drop the rest, whereas the elastic net's ridge part
shares the weight across the group and keeps them in or out together (the ``grouping effect'').

\medskip
\textbf{(d) Collinearity and VIF.}
Collinearity does not bias the OLS estimates, but it makes them \emph{imprecise and unstable}:
it inflates their variances and standard errors, widens confidence intervals, and can make
individual coefficients change sign or lose significance even when the predictors are jointly
important (predictions may still be fine). With $R_j^2=0.9$,
\[
  \mathrm{VIF}_j=\frac{1}{1-0.9}=\frac{1}{0.1}=\mathbf{10},
\]
so the variance of $\hat\beta_j$ is $10$ times larger than it would be with an uncorrelated
$X_j$, and its standard error is inflated by $\sqrt{10}=\mathbf{3.162}$. As a rule of thumb a
$\mathrm{VIF}$ above $5$--$10$ signals problematic collinearity.

\medskip
\textbf{(e) Matching.}
\begin{itemize}[nosep,leftmargin=5mm]
  \item (i) \textbf{Linear regression}: quantitative response, goal is \emph{inference}
    (effect size with confidence interval), $n\gg p$ --- exactly what OLS with standard errors
    provides.
  \item (ii) \textbf{Lasso}: $p\gg n$ ($8{,}000$ vs.\ $80$) makes OLS infeasible; the $\ell_1$
    penalty regularises and selects a short list of metabolites.
  \item (iii) \textbf{Logistic regression}: binary response with interpretable coefficients
    ($e^{\hat\beta_j}$ = odds ratios) --- suitable for reporting how each characteristic changes
    the odds.
  \item (iv) \textbf{Random forest}: captures non-linearities and interactions automatically
    from many predictors and is accurate and robust; interpretability is not required, only
    forecast accuracy.
\end{itemize}
\end{solutionbox}
\fi

\vspace{6mm}
\begin{center}
\rule{0.35\textwidth}{0.4pt}\\[1mm]
\textit{End of the examination --- good luck!}
\end{center}

\end{document}
"
% ============================================================
\documentclass[11pt,a4paper]{article}
\usepackage[utf8]{inputenc}
\usepackage[T1]{fontenc}
\usepackage[english]{babel}
\usepackage[margin=2.2cm]{geometry}
\usepackage{amsmath,amssymb}
\usepackage{booktabs,array}
\usepackage{enumitem}
\usepackage{xcolor}
\usepackage{tcolorbox}
\tcbuselibrary{skins,breakable}

\setlength{\parindent}{0pt}

% Teal solution box (only shown when \withsolutions is defined)
\newtcolorbox{solutionbox}{enhanced,breakable,colback=teal!5,colframe=teal!55!black,
  boxrule=0pt,leftrule=2.5pt,arc=2pt,fonttitle=\bfseries,title=Solution,
  left=2mm,right=2mm,top=1mm,bottom=1.5mm}

\newcommand{\problem}[2]{\par\vspace{7mm}\noindent{\large\textbf{Problem #1}}%
  \hfill{\large\textbf{(#2 points)}}\par\nopagebreak\vspace{0.6mm}%
  \noindent\rule{\textwidth}{0.6pt}\par\nopagebreak\vspace{2.5mm}}
\newcommand{\pp}[1]{\textbf{[#1~P]}\enspace}

\begin{document}

% ------------------------------------------------------------
% Header block
% ------------------------------------------------------------
\begin{center}
  {\Large\textbf{HSBI --- Bielefeld University of Applied Sciences and Arts}}\\[1.2mm]
  {\large Quantitative Research Methods}\\[1mm]
  {\normalsize Summer Semester 2026 \quad$\cdot$\quad Examiner: Prof.\ Dr.\ Christoph Weisser}\\[4.5mm]
  {\LARGE\textbf{Final Mock Examination --- Set C (Lectures 1--12)}}\\[1.2mm]
  {\normalsize Comprehensive coverage of ISLP Chapters 2--8, 10 and 13,\\
   with emphasis on Chapters 7 (splines/GAMs), 8 (trees), 10 (deep learning) and 13 (multiple testing)}%
  \ifdefined\withsolutions\\[2.2mm]{\large\textcolor{teal!55!black}{\textbf{--- Version with detailed solutions ---}}}\fi
\end{center}

\vspace{2.5mm}
\begin{center}
\begin{tabular}{@{}ll@{}}
  \textbf{Duration:} & \textbf{120 minutes}\\
  \textbf{Total credit:} & \textbf{120 points}\\
  \textbf{Allowed aids:} & non-programmable calculator; one handwritten A4 sheet of notes (both sides)\\
\end{tabular}
\end{center}

\vspace{3mm}
Name: \rule[-1mm]{0.38\textwidth}{0.4pt}\hfill Matriculation no.: \rule[-1mm]{0.24\textwidth}{0.4pt}

\vspace{4.5mm}
\begin{center}
\renewcommand{\arraystretch}{1.35}
\begin{tabular}{|l|c|c|c|c|c|c|c|c|}
\hline
\textbf{Problem} & \textbf{P1} & \textbf{P2} & \textbf{P3} & \textbf{P4} & \textbf{P5} & \textbf{P6} & \textbf{P7} & \textbf{Total}\\\hline
Maximum points & 16 & 8 & 20 & 16 & 20 & 20 & 20 & \textbf{120}\\\hline
Achieved points & \hspace{8mm} & \hspace{8mm} & \hspace{8mm} & \hspace{8mm} & \hspace{8mm} & \hspace{8mm} & \hspace{8mm} & \hspace{10mm}\\\hline
\end{tabular}
\end{center}

\vspace{2.5mm}
\textbf{Instructions}
\begin{itemize}[nosep,leftmargin=5.5mm]
  \item Justify all answers. Numerical results without a derivation or a brief reasoning receive no credit.
  \item Round final numerical results to 3 decimal places unless stated otherwise. Small deviations caused by carrying rounded intermediate results are accepted.
  \item All numerical constants you need (powers of $e$, logarithms, $(0.95)^{40}$, $t$-quantiles) are provided in the problems.
  \item The number of points per problem roughly equals the recommended working time in minutes.
\end{itemize}
\vspace{1mm}\noindent\rule{\textwidth}{0.6pt}

% ============================================================
\problem{1: Moving beyond linearity --- splines}{16}
{\small\itshape Lecture slides: Chapter 7 --- Moving Beyond Linearity (Lecture 9).}\par\vspace{1mm}

\begin{enumerate}[label=(\alph*),leftmargin=7mm,itemsep=2.5mm]
  \item \pp{4} A \emph{cubic spline} is a piecewise cubic polynomial that is continuous and has
    continuous first and second derivatives at each knot. Derive the number of degrees of freedom
    (free parameters, counting the intercept) of a cubic spline with $K$ interior knots by
    counting the coefficients of the polynomial pieces and subtracting the constraints imposed
    at the knots. Evaluate your formula for $K=3$ interior knots.
  \item \pp{4} How many degrees of freedom does a \emph{natural} cubic spline with the same
    $K=3$ knots use? State the additional constraint a natural spline imposes, say in which
    region of the predictor space it takes effect, and explain why it is useful in practice.
  \item \pp{4} A \emph{smoothing spline} is the function $g$ that minimises
    \[
      \sum_{i=1}^{n}\bigl(y_i-g(x_i)\bigr)^2+\lambda\int g''(t)^2\,dt .
    \]
    Explain the role of each of the two terms, and describe the behaviour of the solution
    $\hat{g}$ in the two limiting cases $\lambda\to 0$ and $\lambda\to\infty$.
  \item \pp{4} Write down the \emph{truncated-power basis} representation of a cubic spline with
    a single knot at $\xi$, define the truncated-power basis function explicitly, and explain in
    one sentence why adding this basis function does not destroy the smoothness of the fit at
    $\xi$.
\end{enumerate}

\ifdefined\withsolutions
\begin{solutionbox}
\textbf{(a) Degrees of freedom of a cubic spline.}
$K$ interior knots split the range of $X$ into $K+1$ intervals, each carrying a cubic
polynomial with $4$ coefficients: $4(K+1)=4K+4$ parameters. At each knot, continuity of
$g$, $g'$ and $g''$ imposes $3$ constraints, i.e.\ $3K$ constraints in total. Hence
\[
  \mathrm{df}=4(K+1)-3K=K+4 .
\]
For $K=3$: $\mathrm{df}=3+4=\mathbf{7}$ free parameters (equivalently: intercept, $x$, $x^2$,
$x^3$ plus one truncated-power basis function per knot $=4+K$).

\medskip
\textbf{(b) Natural cubic spline.}
A natural spline additionally requires the function to be \emph{linear beyond the boundary
knots}, i.e.\ in the two outer regions $X<\xi_1$ and $X>\xi_K$. Forcing the quadratic and cubic
terms to vanish there removes $2$ parameters per boundary region, i.e.\ $4$ in total:
\[
  \mathrm{df}=(K+4)-4=K,\qquad K=3:\ \mathrm{df}=\mathbf{3}.
\]
Usefulness: polynomials and unconstrained splines are notoriously erratic near the boundaries,
where data are sparse; the linearity constraint stabilises the fit there
(\emph{lower variance} of predictions at and beyond the boundary), at the price of a small
amount of additional bias.

\medskip
\textbf{(c) Smoothing spline.}
The first term (RSS) rewards fidelity to the data; the second term penalises roughness, since
$g''$ measures the curvature (wiggliness) of $g$, and $\lambda\ge 0$ is the tuning parameter
that trades the two off.
\emph{$\lambda\to 0$:} the penalty disappears; $\hat{g}$ can be made arbitrarily wiggly and
\emph{interpolates} the training observations exactly (zero training RSS, severe overfitting;
effective df $\to n$).
\emph{$\lambda\to\infty$:} any curvature is infinitely expensive, so $g''\equiv 0$ is enforced;
$\hat{g}$ becomes a straight line and equals the \emph{least-squares regression line}
(effective df $\to 2$). Between the extremes, $\lambda$ (chosen e.g.\ by LOOCV) controls the
effective degrees of freedom, which decrease smoothly from $n$ to $2$.

\medskip
\textbf{(d) Truncated-power basis with one knot $\xi$.}
\[
  f(x)=\beta_0+\beta_1x+\beta_2x^2+\beta_3x^3+\beta_4\,(x-\xi)^3_+,
  \qquad
  (x-\xi)^3_+=\begin{cases}(x-\xi)^3, & x>\xi\\[0.5mm] 0, & x\le\xi\end{cases}
\]
The function $(x-\xi)^3_+$ and its first and second derivatives are all $0$ at $x=\xi$, so
adding it changes only the third derivative at the knot; the fit remains continuous with
continuous first and second derivatives --- exactly the defining smoothness of a cubic spline.
\end{solutionbox}
\fi

% ============================================================
\problem{2: Generalised additive models}{8}
{\small\itshape Lecture slides: Chapter 7 --- Moving Beyond Linearity (Lecture 9).}\par\vspace{1mm}

\begin{enumerate}[label=(\alph*),leftmargin=7mm,itemsep=2.5mm]
  \item \pp{3} Write down the general form of a \emph{generalised additive model} (GAM) for a
    quantitative response $Y$ and predictors $X_1,\dots,X_p$, and explain in one sentence how it
    generalises multiple linear regression.
  \item \pp{3} State one important \emph{advantage} of GAMs concerning the interpretation of the
    fitted functions $f_j$, and the main structural \emph{limitation} of the additive form ---
    including how this limitation can be repaired.
  \item \pp{2} Describe in one or two sentences how the \emph{backfitting} algorithm fits a GAM.
\end{enumerate}

\ifdefined\withsolutions
\begin{solutionbox}
\textbf{(a) Model form.}
\[
  Y=\beta_0+f_1(X_1)+f_2(X_2)+\dots+f_p(X_p)+\varepsilon .
\]
Multiple linear regression is the special case $f_j(X_j)=\beta_jX_j$; a GAM replaces each linear
component by a smooth non-linear function $f_j$ (e.g.\ a natural or smoothing spline, or local
regression), while keeping the \emph{additive} structure.

\medskip
\textbf{(b) Advantage and limitation.}
\emph{Advantage (interpretability):} because the model is additive, the contribution of each
$X_j$ to $Y$ is captured by its own function $f_j$, which can be examined and plotted
individually while the other variables are held fixed --- non-linear effects remain as easy to
read as in a coefficient table, one panel per predictor.
\emph{Limitation:} additivity excludes \emph{interactions}: the effect of $X_j$ cannot depend
on the value of $X_k$. Repair: add interaction terms manually, e.g.\ a product term
$X_j\times X_k$ or a low-dimensional smooth surface $f_{jk}(X_j,X_k)$ as an extra component.

\medskip
\textbf{(c) Backfitting.}
Cycle through the predictors and repeatedly re-fit each $f_j$ (by a smoother, e.g.\ a spline)
to the \emph{partial residuals} $r_i=y_i-\beta_0-\sum_{k\ne j}f_k(x_{ik})$, holding all other
fitted functions fixed; iterate until the $f_j$ stop changing.
\end{solutionbox}
\fi

% ============================================================
\problem{3: Regression and classification trees by hand}{20}
{\small\itshape Lecture slides: Chapter 8 --- Tree-Based Methods (Lecture 10).}\par\vspace{1mm}

A logistics company records for $n=6$ regional depots the number of delivery vans in operation
$x$ and the number of parcels delivered that day $y$ (in 1{,}000 parcels):

\begin{center}
\begin{tabular}{lcccccc}
\toprule
Depot $i$ & 1 & 2 & 3 & 4 & 5 & 6\\
\midrule
Vans $x_i$ & 1 & 2 & 3 & 4 & 5 & 6\\
Parcels $y_i$ & 3 & 5 & 4 & 9 & 11 & 10\\
\bottomrule
\end{tabular}
\end{center}

A regression tree with a \emph{single} split at cutpoint $s$ partitions the data into
$R_1=\{i: x_i<s\}$ and $R_2=\{i: x_i\ge s\}$ and predicts the region mean $\hat{y}_{R_m}$ in
each region.

\begin{enumerate}[label=(\alph*),leftmargin=7mm,itemsep=2.5mm]
  \item \pp{8} For each of the two candidate cutpoints $s=2.5$ and $s=3.5$, determine the two
    regions, the region means and the resulting total residual sum of squares
    $\operatorname{RSS}(s)=\sum_{m=1}^{2}\sum_{i\in R_m}(y_i-\hat{y}_{R_m})^2$.
    Which cutpoint does recursive binary splitting choose, and why?
  \item \pp{2} Using the chosen split, predict the number of parcels for a new depot that will
    operate $x=5$ vans.
  \item \pp{6} A \emph{classification} tree for a different task has a terminal node containing
    $9$ training observations from three product categories: $5$ ``Electronics'', $3$
    ``Clothing'' and $1$ ``Books''. Compute the class proportions, the Gini index
    $G=\sum_{k}\hat{p}_{k}(1-\hat{p}_{k})$ and the entropy $D=-\sum_{k}\hat{p}_{k}\ln\hat{p}_{k}$
    of this node \emph{(use $\ln(5/9)=-0.5878$, $\ln(1/3)=-1.0986$ and $\ln(1/9)=-2.1972$)}.
    Which class does the node predict?
  \item \pp{4} In \emph{cost-complexity pruning} one minimises, for a subtree
    $T\subseteq T_0$ of the large tree $T_0$,
    \[
      \sum_{m=1}^{|T|}\ \sum_{i:\,x_i\in R_m}\bigl(y_i-\hat{y}_{R_m}\bigr)^2+\alpha\,|T| .
    \]
    Explain the role of the tuning parameter $\alpha$: what happens for $\alpha=0$, what happens
    as $\alpha$ increases, and how is $\alpha$ chosen in practice?
\end{enumerate}

\ifdefined\withsolutions
\begin{solutionbox}
\textbf{(a) Evaluating the two cutpoints.}
Root node for reference: $\bar{y}=42/6=7$.

\emph{Cutpoint $s=2.5$:} $R_1=\{1,2\}$ with $y$-values $\{3,5\}$, mean $\hat{y}_{R_1}=4$;
$R_2=\{3,4,5,6\}$ with $y$-values $\{4,9,11,10\}$, mean $\hat{y}_{R_2}=34/4=8.5$.
\begin{align*}
  \operatorname{RSS}_{R_1}&=(3-4)^2+(5-4)^2=2,\\
  \operatorname{RSS}_{R_2}&=(4-8.5)^2+(9-8.5)^2+(11-8.5)^2+(10-8.5)^2=20.25+0.25+6.25+2.25=29,
\end{align*}
so $\operatorname{RSS}(2.5)=2+29=\mathbf{31}$.

\emph{Cutpoint $s=3.5$:} $R_1=\{1,2,3\}$ with $y$-values $\{3,5,4\}$, mean
$\hat{y}_{R_1}=12/3=4$; $R_2=\{4,5,6\}$ with $y$-values $\{9,11,10\}$, mean
$\hat{y}_{R_2}=30/3=10$.
\[
  \operatorname{RSS}_{R_1}=(3-4)^2+(5-4)^2+(4-4)^2=2,\qquad
  \operatorname{RSS}_{R_2}=(9-10)^2+(11-10)^2+(10-10)^2=2,
\]
so $\operatorname{RSS}(3.5)=2+2=\mathbf{4}$.

Since $4<31$, recursive binary splitting is \emph{greedy} and picks the cutpoint with the
smallest RSS: $\mathbf{s=3.5}$. (It cleanly separates the three low-volume depots, mean $4$,
from the three high-volume depots, mean $10$; in fact $s=3.5$ is the best of all five possible
cutpoints.)

\medskip
\textbf{(b) Prediction for $x=5$.}
$x=5\ge 3.5$, so the new depot falls into $R_2$ and the tree predicts
$\hat{y}=\hat{y}_{R_2}=\mathbf{10}$ (i.e.\ $10{,}000$ parcels).

\medskip
\textbf{(c) Node impurity (three classes).}
$\hat{p}_{\text{E}}=5/9$, $\hat{p}_{\text{C}}=3/9=1/3$, $\hat{p}_{\text{B}}=1/9$.
\[
  G=\tfrac{5}{9}\cdot\tfrac{4}{9}+\tfrac{3}{9}\cdot\tfrac{6}{9}+\tfrac{1}{9}\cdot\tfrac{8}{9}
   =\tfrac{20+18+8}{81}=\tfrac{46}{81}=\mathbf{0.568}.
\]
\[
  D=-\Bigl(\tfrac{5}{9}\ln\tfrac{5}{9}+\tfrac{1}{3}\ln\tfrac{1}{3}+\tfrac{1}{9}\ln\tfrac{1}{9}\Bigr)
   =-\bigl(\tfrac{5}{9}(-0.5878)+\tfrac{1}{3}(-1.0986)+\tfrac{1}{9}(-2.1972)\bigr)
\]
\[
   =0.3266+0.3662+0.2441=\mathbf{0.937}.
\]
The node predicts the majority class, \textbf{Electronics} (with estimated probability
$5/9\approx 0.556$). Both $G$ and $D$ would be $0$ for a pure node; here the node is quite
impure because no class dominates strongly.

\medskip
\textbf{(d) Role of $\alpha$.}
$\alpha$ is the price per terminal node: it trades off the fit of the subtree (first term)
against its complexity $|T|$.
For $\alpha=0$ the criterion equals the training RSS and the full tree $T_0$ is optimal.
As $\alpha$ increases, additional leaves must ``pay for themselves''; branches whose RSS
reduction is smaller than $\alpha$ are pruned, so the optimal subtree becomes smaller ---
one obtains a \emph{nested sequence} of subtrees as $\alpha$ grows (weakest-link pruning).
In practice $\alpha$ is selected by \emph{cross-validation} (choose the $\alpha$ with the
smallest CV error, then refit on all data), which balances bias against variance instead of
optimising training fit.
\end{solutionbox}
\fi

% ============================================================
\problem{4: Ensembles --- bagging, random forests, boosting}{16}
{\small\itshape Lecture slides: Chapter 8 --- Tree-Based Methods (Lecture 10).}\par\vspace{1mm}

\begin{enumerate}[label=(\alph*),leftmargin=7mm,itemsep=2.5mm]
  \item \pp{6} Let $\hat{f}_1(x),\dots,\hat{f}_B(x)$ be $B$ identically distributed (but not
    independent) tree predictions at a fixed point $x$, each with variance $\sigma^2$ and with
    pairwise correlation $\rho$. The variance of the bagged average is
    \[
      \operatorname{Var}\Bigl(\tfrac{1}{B}\sum_{b=1}^{B}\hat{f}_b(x)\Bigr)
      =\rho\,\sigma^{2}+\frac{1-\rho}{B}\,\sigma^{2}.
    \]
    (i) Explain briefly why this formula shows that bagging reduces variance, and which term
    limits the reduction. (ii) Evaluate the variance for $B=50$, $\rho=0.4$ and $\sigma^2=16$,
    and compare it with the variance of a single tree. (iii) What value does the variance
    approach as $B\to\infty$?
  \item \pp{4} How do \emph{random forests} modify bagging at each split, and what is the
    recommended subset size $m$ for \emph{regression} as a function of $p$ (evaluate for
    $p=30$)? Explain, with reference to the formula in (a), why this modification improves on
    bagging.
  \item \pp{4} Name the three tuning parameters of \emph{boosting} for regression trees and
    state in one sentence each what they control. What is the ``slow learning'' idea behind
    boosting?
  \item \pp{2} What does the \emph{out-of-bag} (OOB) error of a bagged model estimate, and why
    is it available essentially for free?
\end{enumerate}

\ifdefined\withsolutions
\begin{solutionbox}
\textbf{(a) Variance of the bagged average.}
(i) The second term $\frac{1-\rho}{B}\sigma^2$ is driven to $0$ by averaging over many trees
($B$ large), so averaging removes the ``independent part'' of the trees' variability ---
this is the variance reduction of bagging. The first term $\rho\sigma^2$ does \emph{not}
depend on $B$: because the bootstrap trees are grown on overlapping samples of the same data
(and share the same strong predictors), they are positively correlated, and this correlated
part of the variance cannot be averaged away. It limits the reduction.
(ii) For $B=50$, $\rho=0.4$, $\sigma^2=16$:
\[
  \operatorname{Var}=0.4\cdot 16+\frac{1-0.4}{50}\cdot 16
  =6.4+\frac{0.6\cdot 16}{50}=6.4+0.192=\mathbf{6.592},
\]
compared with $\sigma^2=16$ for a single tree --- a reduction of about $58.8\%$.
(iii) As $B\to\infty$ the variance approaches the floor $\rho\sigma^2=0.4\cdot 16=\mathbf{6.4}$.
(Increasing $B$ further does not overfit; it only stabilises the average.)

\medskip
\textbf{(b) Random forests.}
At each split, only a \emph{random subset} of $m$ of the $p$ predictors is considered as
splitting candidates; a fresh subset is drawn at every split. Recommended for regression:
$m\approx p/3$, so for $p=30$: $m=30/3=\mathbf{10}$.
Why it helps: with bagging, a dominant predictor is used in the top split of almost every tree,
so the trees look alike and $\rho$ is high. Excluding it from $\approx (p-m)/p$ of the split
searches \emph{decorrelates} the trees, i.e.\ lowers $\rho$ --- and by (a) this lowers the
irreducible floor $\rho\sigma^2$ of the ensemble variance, which mere averaging cannot touch.

\medskip
\textbf{(c) Boosting.}
\begin{itemize}[nosep,leftmargin=5mm]
  \item \emph{Number of trees $B$:} how many small trees are added sequentially; unlike
    bagging, boosting \emph{can} overfit if $B$ is too large, so $B$ is chosen by
    cross-validation.
  \item \emph{Shrinkage $\lambda$} (learning rate, e.g.\ $0.01$ or $0.001$): scales the
    contribution of each new tree and thus the speed of learning; smaller $\lambda$ usually
    needs larger $B$.
  \item \emph{Number of splits $d$ per tree} (interaction depth): the complexity of each tree
    and the interaction order of the model; often even $d=1$ (stumps, additive model) works
    well.
\end{itemize}
\emph{Slow learning:} each tree is fitted to the current \emph{residuals}, and only a small
fraction $\lambda$ of it is added to the model; the model thus improves slowly and exactly in
the regions where it still performs poorly, instead of fitting the data hard in one step ---
which tends to generalise better.

\medskip
\textbf{(d) OOB error.}
Each bootstrap sample leaves out on average about $1/3$ of the observations
(``out-of-bag''). Predicting each observation only with the trees that did \emph{not} use it,
and aggregating these predictions, yields the OOB error --- a valid estimate of the
\emph{test error} of the bagged model. It is free because it reuses the bootstrap side
product; no separate validation set or cross-validation loop is required.
\end{solutionbox}
\fi

% ============================================================
\problem{5: Neural networks by hand}{20}
{\small\itshape Lecture slides: Chapter 10 --- Deep Learning (Lecture 11).}\par\vspace{1mm}

Consider a tiny feed-forward network for a quantitative response with input
$x=(x_1,x_2)=(2,3)$, one hidden layer with two ReLU units ($g(z)=\max\{0,z\}$) and a linear
output unit:
\[
  A_1=g\bigl(-3+2\,x_1+1\,x_2\bigr),\qquad
  A_2=g\bigl(2-2\,x_1-1\,x_2\bigr),\qquad
  f(x)=-2+3\,A_1+2\,A_2 .
\]

\begin{enumerate}[label=(\alph*),leftmargin=7mm,itemsep=2.5mm]
  \item \pp{8} Perform the forward pass step by step: compute both pre-activations, both
    activations $A_1,A_2$ and the network output $f(x)$. Comment briefly on what the ReLU did
    to the second unit and why such non-linear activations are essential.
  \item \pp{6} A classification network with $K=3$ classes produces the logits (output-layer
    pre-activations) $z=(z_1,z_2,z_3)=(2.0,\ 0.0,\ -1.0)$ for one observation whose true class
    is class~2. Compute all three softmax probabilities and the cross-entropy loss
    $-\ln \hat{p}_{\text{true}}$.
    \emph{(Use $e^{2}=7.389$, $e^{0}=1.000$, $e^{-1}=0.368$ and $\ln 8.757=2.170$.)}
  \item \pp{4} Count the parameters of a fully connected network with $p=6$ inputs, one hidden
    layer with $k=5$ units and an output layer with $K=3$ units, where every hidden and output
    unit has a bias. Show the count separately for the two weight layers.
  \item \pp{2} A convolutional layer receives a $64\times 64$ image and uses filters of size
    $F=3$ with padding $P=1$ and stride $S=2$. Compute the spatial size of the output feature
    map from $\bigl\lfloor(W-F+2P)/S\bigr\rfloor+1$ (note the stride is $2$, so the division is
    rounded down).
\end{enumerate}

\ifdefined\withsolutions
\begin{solutionbox}
\textbf{(a) Forward pass.}
Pre-activations:
\[
  z_1=-3+2\cdot 2+1\cdot 3=-3+4+3=4,\qquad
  z_2=2-2\cdot 2-1\cdot 3=2-4-3=-5 .
\]
ReLU activations:
\[
  A_1=g(4)=\max\{0,4\}=4,\qquad A_2=g(-5)=\max\{0,-5\}=0 .
\]
Output:
\[
  f(x)=-2+3\cdot 4+2\cdot 0=-2+12+0=\mathbf{10}.
\]
The ReLU \emph{switched the second unit off}: its negative pre-activation is clipped to $0$,
so unit 2 contributes nothing for this input (but would for other inputs). This thresholding
is the source of non-linearity: without a non-linear $g$, the network would collapse into a
single linear function of $x$, no matter how many layers it has.

\medskip
\textbf{(b) Softmax and cross-entropy.}
Exponentials: $e^{2.0}=7.389$, $e^{0.0}=1.000$, $e^{-1.0}=0.368$; their sum is
$7.389+1.000+0.368=8.757$. Softmax probabilities:
\[
  \hat{p}_1=\frac{7.389}{8.757}=0.844,\qquad
  \hat{p}_2=\frac{1.000}{8.757}=0.114,\qquad
  \hat{p}_3=\frac{0.368}{8.757}=0.042,
\]
which sum to $1.000$. Cross-entropy loss for the true class $2$:
\[
  -\ln\hat{p}_2=-\ln\frac{1.000}{8.757}=-\bigl(0-\ln 8.757\bigr)=\ln 8.757=\mathbf{2.170}.
\]
(Note that the model's most probable class is class~1, but the loss is evaluated at the
\emph{true} class~2, which received only probability $0.114$ --- hence the sizeable loss.)

\medskip
\textbf{(c) Parameter count for $6\to 5\to 3$.}
Input $\to$ hidden: $6\cdot 5=30$ weights $+\,5$ biases $=35$.
Hidden $\to$ output: $5\cdot 3=15$ weights $+\,3$ biases $=18$.
Total: $35+18=\mathbf{53}$ parameters.

\medskip
\textbf{(d) Convolution output size.}
\[
  \Bigl\lfloor\frac{W-F+2P}{S}\Bigr\rfloor+1
  =\Bigl\lfloor\frac{64-3+2\cdot 1}{2}\Bigr\rfloor+1
  =\Bigl\lfloor\frac{63}{2}\Bigr\rfloor+1=\lfloor 31.5\rfloor+1=31+1=32,
\]
so the output feature map is $\mathbf{32\times 32}$ --- the stride $S=2$ roughly halves the
$64\times 64$ input (the floor discards the fractional part before adding $1$).
\end{solutionbox}
\fi

% ============================================================
\problem{6: Multiple testing}{20}
{\small\itshape Lecture slides: Chapter 13 --- Multiple Testing (Lecture 12).}\par\vspace{1mm}

\begin{enumerate}[label=(\alph*),leftmargin=7mm,itemsep=2.5mm]
  \item \pp{4} A lab performs $m=40$ \emph{independent} hypothesis tests, each at level
    $\alpha=0.05$, and suppose all $40$ null hypotheses are true. Define the family-wise error
    rate (FWER) and compute it for this situation \emph{(use $(0.95)^{40}=0.1285$)}. Interpret
    the result in one sentence.
  \item \pp{7} A research group screens $m=6$ candidate biomarkers; the ordered $p$-values are
    \begin{center}
    \begin{tabular}{lcccccc}
    \toprule
    $j$ & 1 & 2 & 3 & 4 & 5 & 6\\
    \midrule
    $p_{(j)}$ & 0.001 & 0.008 & 0.011 & 0.030 & 0.041 & 0.260\\
    \bottomrule
    \end{tabular}
    \end{center}
    Control the FWER at $\alpha=0.05$ (i) with \emph{Bonferroni} and (ii) with the
    \emph{Holm step-down} procedure. Give the relevant thresholds, state for each method
    exactly which hypotheses are rejected, and explain why Holm can never reject fewer
    hypotheses than Bonferroni.
  \item \pp{6} Apply the \emph{Benjamini--Hochberg} (BH) procedure with FDR level $q=0.05$ to
    the same six $p$-values: give the thresholds $jq/m$, determine which hypotheses are
    rejected, and state precisely what ``FDR control at $q=0.05$'' guarantees for the
    resulting set of rejections.
  \item \pp{3} A pharmaceutical company runs a \emph{confirmatory} phase-III trial that tests
    $m=5$ pre-specified endpoints; a single false-positive endpoint could lead to the approval
    of an ineffective drug. Should they control the FWER or the FDR? Justify your
    recommendation.
\end{enumerate}

\ifdefined\withsolutions
\begin{solutionbox}
\textbf{(a) FWER for $m=40$ independent tests.}
$\text{FWER}=\Pr(\text{at least one Type I error})=\Pr(V\ge 1)$, where $V$ is the number of
true nulls that are (falsely) rejected. With independence and all nulls true:
\[
  \text{FWER}=1-\Pr(\text{no false rejection})=1-(1-0.05)^{40}=1-(0.95)^{40}
  =1-0.1285=\mathbf{0.8715}.
\]
Interpretation: even though each single test is performed at the $5\%$ level, there is about an
$87\%$ chance of at least one false discovery among the $40$ tests --- testing many hypotheses
without adjustment almost guarantees spurious findings.

\medskip
\textbf{(b) Bonferroni and Holm at $\alpha=0.05$.}
\emph{Bonferroni:} reject $H_{(j)}$ iff $p_{(j)}\le\alpha/m=0.05/6=0.00833$.
$p_{(1)}=0.001\le 0.00833$ \checkmark, $p_{(2)}=0.008\le 0.00833$ \checkmark,
$p_{(3)}=0.011>0.00833$ $\times$.
$\Rightarrow$ Bonferroni rejects $H_{(1)}$ and $H_{(2)}$: \textbf{2 rejections}.

\emph{Holm:} compare $p_{(j)}$ with $\alpha/(m-j+1)$ step by step and stop at the first failure:
\begin{center}
\begin{tabular}{lcccccc}
\toprule
$j$ & 1 & 2 & 3 & 4 & 5 & 6\\
\midrule
$p_{(j)}$ & 0.001 & 0.008 & 0.011 & 0.030 & 0.041 & 0.260\\
$\alpha/(m-j+1)$ & 0.00833 & 0.0100 & 0.0125 & 0.0167 & 0.0250 & 0.0500\\
reject? & \checkmark & \checkmark & \checkmark & stop & --- & ---\\
\bottomrule
\end{tabular}
\end{center}
$0.001\le 0.00833$, $0.008\le 0.0100$, $0.011\le 0.0125$, but $0.030>0.0167$: stop.
$\Rightarrow$ Holm rejects $H_{(1)},H_{(2)},H_{(3)}$: \textbf{3 rejections}.

\emph{Why Holm $\supseteq$ Bonferroni:} Holm's thresholds $\alpha/(m-j+1)$ are $\ge\alpha/m$
for every $j$ (with equality only at $j=1$), so every $p$-value that passes Bonferroni also
passes Holm's less stringent later hurdles --- Holm is uniformly more powerful, yet still
controls the FWER at $\alpha$ without any independence assumption. Here Holm rejects one more
hypothesis ($H_{(3)}$) than Bonferroni.

\medskip
\textbf{(c) Benjamini--Hochberg at $q=0.05$.}
Thresholds $jq/m=j\cdot 0.05/6$:
\begin{center}
\begin{tabular}{lcccccc}
\toprule
$j$ & 1 & 2 & 3 & 4 & 5 & 6\\
\midrule
$p_{(j)}$ & 0.001 & 0.008 & 0.011 & 0.030 & 0.041 & 0.260\\
$jq/m$ & 0.00833 & 0.0167 & 0.0250 & 0.0333 & 0.0417 & 0.0500\\
$p_{(j)}\le jq/m$? & yes & yes & yes & yes & \textbf{yes} & no\\
\bottomrule
\end{tabular}
\end{center}
$L=$ largest $j$ with $p_{(j)}\le jq/m$: $j=5$ (since $0.041\le 0.0417$, while
$0.260>0.0500$). BH rejects \emph{all} $H_{(j)}$ with $j\le L$: $H_{(1)},\dots,H_{(5)}$:
\textbf{5 rejections}.
Guarantee: $\text{FDR}=E\bigl[V/\max(R,1)\bigr]\le q=0.05$ --- \emph{on average} at most
$5\%$ of the rejected hypotheses are false discoveries. It does \emph{not} guarantee that at
most $5\%$ of these particular $5$ rejections are false, nor does it bound the probability of
making any false rejection (that would be FWER control). Note BH is the most liberal of the
three here ($5$ rejections vs.\ $3$ for Holm and $2$ for Bonferroni).

\medskip
\textbf{(d) Confirmatory phase-III trial.}
Control the \textbf{FWER}. This is a confirmatory study with only $m=5$ pre-specified
endpoints, and the cost of even a single false positive is severe (approval of an ineffective
drug, patient harm, regulatory failure). FWER control (e.g.\ Holm at $\alpha=0.05$) bounds the
probability of \emph{any} false rejection at $5\%$, which is exactly the guarantee a regulator
demands. FDR control would only limit the \emph{expected proportion} of false endpoints and
would tolerate an occasional false positive --- acceptable in large exploratory screens, but
not when each individual claim must stand on its own.
\end{solutionbox}
\fi

% ============================================================
\problem{7: Cumulative essentials (Chapters 2--6)}{20}
{\small\itshape Lecture slides: Chapters 2--6 (Lectures 1--8).}\par\vspace{1mm}

\begin{enumerate}[label=(\alph*),leftmargin=7mm,itemsep=2.5mm]
  \item \pp{4} The \emph{bootstrap} draws a resample of size $n$ with replacement from the $n$
    observations. (i) Show that the probability that a given observation is \emph{not} selected
    into a particular bootstrap sample equals $(1-1/n)^n$, and give its limit as $n\to\infty$
    \emph{(use $e^{-1}=0.368$)}. (ii) Hence what fraction of the original observations appears
    (at least once) in a typical bootstrap sample as $n\to\infty$? (iii) Compute this inclusion
    probability exactly for $n=5$ \emph{(use $0.8^{5}=0.328$)}.
  \item \pp{4} At a fixed test point $x_0$ the expected test error of an estimate $\hat f$
    decomposes as
    $E\bigl[(y_0-\hat f(x_0))^2\bigr]
      =\operatorname{Var}(\hat f(x_0))+\bigl[\operatorname{Bias}(\hat f(x_0))\bigr]^2
       +\operatorname{Var}(\varepsilon)$.
    Two candidate methods give, at $x_0$:
    \begin{center}
    \begin{tabular}{lccc}
    \toprule
    Method & $\operatorname{Var}(\hat f)$ & $\operatorname{Bias}(\hat f)$ & $\operatorname{Var}(\varepsilon)$\\
    \midrule
    (1) flexible spline & 3.5 & $-1.0$ & 2.0\\
    (2) linear model & 0.5 & $-2.5$ & 2.0\\
    \bottomrule
    \end{tabular}
    \end{center}
    Compute the expected test MSE of each method, say which you would deploy at $x_0$, and state
    the lowest test MSE \emph{any} method could achieve here.
  \item \pp{3} Write down the \emph{elastic net} penalty and explain what it inherits from the
    lasso and what it inherits from ridge regression. In which situation does the elastic net
    clearly outperform the lasso?
  \item \pp{3} Two predictors in a regression have correlation $0.95$. Explain what such
    \emph{collinearity} does to the least-squares coefficient estimates. The variance inflation
    factor is $\mathrm{VIF}_j=1/(1-R_j^2)$, where $R_j^2$ is from regressing $X_j$ on the other
    predictors; compute and interpret $\mathrm{VIF}_j$ for $R_j^2=0.9$ \emph{(use
    $\sqrt{10}=3.162$)}.
  \item \pp{6} Match each scenario to exactly one method from
    \{\emph{linear regression}, \emph{lasso}, \emph{logistic regression},
    \emph{random forest}\} (each method used exactly once) and justify each choice in one
    sentence:
    \begin{enumerate}[label=(\roman*),nosep,topsep=1mm]
      \item Estimate, with a $95\%$ confidence interval, how much one additional year of
        schooling raises hourly wage, from $n=3{,}000$ workers and a handful of predictors.
      \item From $n=80$ patients with $p=8{,}000$ metabolite measurements, predict a continuous
        biomarker level and identify a short list of relevant metabolites.
      \item Predict whether a customer will churn (yes/no) and report to management how each
        customer characteristic changes the \emph{odds} of churning.
      \item Forecast next-day electricity demand (continuous) from hundreds of weather and
        calendar variables with strong non-linear effects and interactions; only forecast
        accuracy matters.
    \end{enumerate}
\end{enumerate}

\ifdefined\withsolutions
\begin{solutionbox}
\textbf{(a) Bootstrap inclusion probability.}
(i) In each of the $n$ independent draws, a fixed observation is \emph{missed} with probability
$1-\tfrac1n$; since the draws are independent, all $n$ draws miss it with probability
$(1-\tfrac1n)^n$. As $n\to\infty$, $(1-\tfrac1n)^n\to e^{-1}=\mathbf{0.368}$.
(ii) The observation is therefore \emph{included} with probability $1-e^{-1}=1-0.368=0.632$, so
on average about $\mathbf{63.2\%}$ of the original observations appear in a bootstrap sample
(the remaining $\approx 1/3$ are the ``out-of-bag'' observations).
(iii) For $n=5$: $\Pr(\text{included})=1-(1-\tfrac15)^5=1-0.8^{5}=1-0.328=\mathbf{0.672}$.

\medskip
\textbf{(b) Bias--variance decomposition.}
\[
  \text{MSE}_1=\operatorname{Var}+\operatorname{Bias}^2+\operatorname{Var}(\varepsilon)
   =3.5+(-1.0)^2+2.0=3.5+1.0+2.0=\mathbf{6.5},
\]
\[
  \text{MSE}_2=0.5+(-2.5)^2+2.0=0.5+6.25+2.0=\mathbf{8.75}.
\]
Deploy the \textbf{flexible spline} (method 1): its higher estimation variance ($3.5$ vs.\
$0.5$) is more than offset by its much smaller squared bias ($1.0$ vs.\ $6.25$), giving the
lower expected test error ($6.5<8.75$). The lowest test MSE any method can achieve at $x_0$ is
the \emph{irreducible error} $\operatorname{Var}(\varepsilon)=\mathbf{2.0}$; no estimator can go
below it.

\medskip
\textbf{(c) Elastic net.}
The elastic net minimises, for $0\le\alpha\le 1$,
\[
  \sum_{i=1}^n\bigl(y_i-\beta_0-\textstyle\sum_j\beta_jx_{ij}\bigr)^2
  +\lambda\Bigl[\tfrac{1-\alpha}{2}\textstyle\sum_j\beta_j^2+\alpha\textstyle\sum_j|\beta_j|\Bigr],
\]
i.e.\ a weighted mix of the ridge ($\ell_2$) and lasso ($\ell_1$) penalties. From the
\emph{lasso} it inherits the ability to set coefficients \emph{exactly to zero} (sparsity /
variable selection); from \emph{ridge} it inherits the graceful handling of correlated
predictors and the ability to select more than $n$ variables when $p>n$. It clearly beats the
lasso when there are \emph{groups of highly correlated predictors}: the lasso tends to pick one
member of a group essentially at random and drop the rest, whereas the elastic net's ridge part
shares the weight across the group and keeps them in or out together (the ``grouping effect'').

\medskip
\textbf{(d) Collinearity and VIF.}
Collinearity does not bias the OLS estimates, but it makes them \emph{imprecise and unstable}:
it inflates their variances and standard errors, widens confidence intervals, and can make
individual coefficients change sign or lose significance even when the predictors are jointly
important (predictions may still be fine). With $R_j^2=0.9$,
\[
  \mathrm{VIF}_j=\frac{1}{1-0.9}=\frac{1}{0.1}=\mathbf{10},
\]
so the variance of $\hat\beta_j$ is $10$ times larger than it would be with an uncorrelated
$X_j$, and its standard error is inflated by $\sqrt{10}=\mathbf{3.162}$. As a rule of thumb a
$\mathrm{VIF}$ above $5$--$10$ signals problematic collinearity.

\medskip
\textbf{(e) Matching.}
\begin{itemize}[nosep,leftmargin=5mm]
  \item (i) \textbf{Linear regression}: quantitative response, goal is \emph{inference}
    (effect size with confidence interval), $n\gg p$ --- exactly what OLS with standard errors
    provides.
  \item (ii) \textbf{Lasso}: $p\gg n$ ($8{,}000$ vs.\ $80$) makes OLS infeasible; the $\ell_1$
    penalty regularises and selects a short list of metabolites.
  \item (iii) \textbf{Logistic regression}: binary response with interpretable coefficients
    ($e^{\hat\beta_j}$ = odds ratios) --- suitable for reporting how each characteristic changes
    the odds.
  \item (iv) \textbf{Random forest}: captures non-linearities and interactions automatically
    from many predictors and is accurate and robust; interpretability is not required, only
    forecast accuracy.
\end{itemize}
\end{solutionbox}
\fi

\vspace{6mm}
\begin{center}
\rule{0.35\textwidth}{0.4pt}\\[1mm]
\textit{End of the examination --- good luck!}
\end{center}

\end{document}
"
% ============================================================
\documentclass[11pt,a4paper]{article}
\usepackage[utf8]{inputenc}
\usepackage[T1]{fontenc}
\usepackage[english]{babel}
\usepackage[margin=2.2cm]{geometry}
\usepackage{amsmath,amssymb}
\usepackage{booktabs,array}
\usepackage{enumitem}
\usepackage{xcolor}
\usepackage{tcolorbox}
\tcbuselibrary{skins,breakable}

\setlength{\parindent}{0pt}

% Teal solution box (only shown when \withsolutions is defined)
\newtcolorbox{solutionbox}{enhanced,breakable,colback=teal!5,colframe=teal!55!black,
  boxrule=0pt,leftrule=2.5pt,arc=2pt,fonttitle=\bfseries,title=Solution,
  left=2mm,right=2mm,top=1mm,bottom=1.5mm}

\newcommand{\problem}[2]{\par\vspace{7mm}\noindent{\large\textbf{Problem #1}}%
  \hfill{\large\textbf{(#2 points)}}\par\nopagebreak\vspace{0.6mm}%
  \noindent\rule{\textwidth}{0.6pt}\par\nopagebreak\vspace{2.5mm}}
\newcommand{\pp}[1]{\textbf{[#1~P]}\enspace}

\begin{document}

% ------------------------------------------------------------
% Header block
% ------------------------------------------------------------
\begin{center}
  {\Large\textbf{HSBI --- Bielefeld University of Applied Sciences and Arts}}\\[1.2mm]
  {\large Quantitative Research Methods}\\[1mm]
  {\normalsize Summer Semester 2026 \quad$\cdot$\quad Examiner: Prof.\ Dr.\ Christoph Weisser}\\[4.5mm]
  {\LARGE\textbf{Final Mock Examination --- Set C (Lectures 1--12)}}\\[1.2mm]
  {\normalsize Comprehensive coverage of ISLP Chapters 2--8, 10 and 13,\\
   with emphasis on Chapters 7 (splines/GAMs), 8 (trees), 10 (deep learning) and 13 (multiple testing)}%
  \ifdefined\withsolutions\\[2.2mm]{\large\textcolor{teal!55!black}{\textbf{--- Version with detailed solutions ---}}}\fi
\end{center}

\vspace{2.5mm}
\begin{center}
\begin{tabular}{@{}ll@{}}
  \textbf{Duration:} & \textbf{120 minutes}\\
  \textbf{Total credit:} & \textbf{120 points}\\
  \textbf{Allowed aids:} & non-programmable calculator; one handwritten A4 sheet of notes (both sides)\\
\end{tabular}
\end{center}

\vspace{3mm}
Name: \rule[-1mm]{0.38\textwidth}{0.4pt}\hfill Matriculation no.: \rule[-1mm]{0.24\textwidth}{0.4pt}

\vspace{4.5mm}
\begin{center}
\renewcommand{\arraystretch}{1.35}
\begin{tabular}{|l|c|c|c|c|c|c|c|c|}
\hline
\textbf{Problem} & \textbf{P1} & \textbf{P2} & \textbf{P3} & \textbf{P4} & \textbf{P5} & \textbf{P6} & \textbf{P7} & \textbf{Total}\\\hline
Maximum points & 16 & 8 & 20 & 16 & 20 & 20 & 20 & \textbf{120}\\\hline
Achieved points & \hspace{8mm} & \hspace{8mm} & \hspace{8mm} & \hspace{8mm} & \hspace{8mm} & \hspace{8mm} & \hspace{8mm} & \hspace{10mm}\\\hline
\end{tabular}
\end{center}

\vspace{2.5mm}
\textbf{Instructions}
\begin{itemize}[nosep,leftmargin=5.5mm]
  \item Justify all answers. Numerical results without a derivation or a brief reasoning receive no credit.
  \item Round final numerical results to 3 decimal places unless stated otherwise. Small deviations caused by carrying rounded intermediate results are accepted.
  \item All numerical constants you need (powers of $e$, logarithms, $(0.95)^{40}$, $t$-quantiles) are provided in the problems.
  \item The number of points per problem roughly equals the recommended working time in minutes.
\end{itemize}
\vspace{1mm}\noindent\rule{\textwidth}{0.6pt}

% ============================================================
\problem{1: Moving beyond linearity --- splines}{16}
{\small\itshape Lecture slides: Chapter 7 --- Moving Beyond Linearity (Lecture 9).}\par\vspace{1mm}

\begin{enumerate}[label=(\alph*),leftmargin=7mm,itemsep=2.5mm]
  \item \pp{4} A \emph{cubic spline} is a piecewise cubic polynomial that is continuous and has
    continuous first and second derivatives at each knot. Derive the number of degrees of freedom
    (free parameters, counting the intercept) of a cubic spline with $K$ interior knots by
    counting the coefficients of the polynomial pieces and subtracting the constraints imposed
    at the knots. Evaluate your formula for $K=3$ interior knots.
  \item \pp{4} How many degrees of freedom does a \emph{natural} cubic spline with the same
    $K=3$ knots use? State the additional constraint a natural spline imposes, say in which
    region of the predictor space it takes effect, and explain why it is useful in practice.
  \item \pp{4} A \emph{smoothing spline} is the function $g$ that minimises
    \[
      \sum_{i=1}^{n}\bigl(y_i-g(x_i)\bigr)^2+\lambda\int g''(t)^2\,dt .
    \]
    Explain the role of each of the two terms, and describe the behaviour of the solution
    $\hat{g}$ in the two limiting cases $\lambda\to 0$ and $\lambda\to\infty$.
  \item \pp{4} Write down the \emph{truncated-power basis} representation of a cubic spline with
    a single knot at $\xi$, define the truncated-power basis function explicitly, and explain in
    one sentence why adding this basis function does not destroy the smoothness of the fit at
    $\xi$.
\end{enumerate}

\ifdefined\withsolutions
\begin{solutionbox}
\textbf{(a) Degrees of freedom of a cubic spline.}
$K$ interior knots split the range of $X$ into $K+1$ intervals, each carrying a cubic
polynomial with $4$ coefficients: $4(K+1)=4K+4$ parameters. At each knot, continuity of
$g$, $g'$ and $g''$ imposes $3$ constraints, i.e.\ $3K$ constraints in total. Hence
\[
  \mathrm{df}=4(K+1)-3K=K+4 .
\]
For $K=3$: $\mathrm{df}=3+4=\mathbf{7}$ free parameters (equivalently: intercept, $x$, $x^2$,
$x^3$ plus one truncated-power basis function per knot $=4+K$).

\medskip
\textbf{(b) Natural cubic spline.}
A natural spline additionally requires the function to be \emph{linear beyond the boundary
knots}, i.e.\ in the two outer regions $X<\xi_1$ and $X>\xi_K$. Forcing the quadratic and cubic
terms to vanish there removes $2$ parameters per boundary region, i.e.\ $4$ in total:
\[
  \mathrm{df}=(K+4)-4=K,\qquad K=3:\ \mathrm{df}=\mathbf{3}.
\]
Usefulness: polynomials and unconstrained splines are notoriously erratic near the boundaries,
where data are sparse; the linearity constraint stabilises the fit there
(\emph{lower variance} of predictions at and beyond the boundary), at the price of a small
amount of additional bias.

\medskip
\textbf{(c) Smoothing spline.}
The first term (RSS) rewards fidelity to the data; the second term penalises roughness, since
$g''$ measures the curvature (wiggliness) of $g$, and $\lambda\ge 0$ is the tuning parameter
that trades the two off.
\emph{$\lambda\to 0$:} the penalty disappears; $\hat{g}$ can be made arbitrarily wiggly and
\emph{interpolates} the training observations exactly (zero training RSS, severe overfitting;
effective df $\to n$).
\emph{$\lambda\to\infty$:} any curvature is infinitely expensive, so $g''\equiv 0$ is enforced;
$\hat{g}$ becomes a straight line and equals the \emph{least-squares regression line}
(effective df $\to 2$). Between the extremes, $\lambda$ (chosen e.g.\ by LOOCV) controls the
effective degrees of freedom, which decrease smoothly from $n$ to $2$.

\medskip
\textbf{(d) Truncated-power basis with one knot $\xi$.}
\[
  f(x)=\beta_0+\beta_1x+\beta_2x^2+\beta_3x^3+\beta_4\,(x-\xi)^3_+,
  \qquad
  (x-\xi)^3_+=\begin{cases}(x-\xi)^3, & x>\xi\\[0.5mm] 0, & x\le\xi\end{cases}
\]
The function $(x-\xi)^3_+$ and its first and second derivatives are all $0$ at $x=\xi$, so
adding it changes only the third derivative at the knot; the fit remains continuous with
continuous first and second derivatives --- exactly the defining smoothness of a cubic spline.
\end{solutionbox}
\fi

% ============================================================
\problem{2: Generalised additive models}{8}
{\small\itshape Lecture slides: Chapter 7 --- Moving Beyond Linearity (Lecture 9).}\par\vspace{1mm}

\begin{enumerate}[label=(\alph*),leftmargin=7mm,itemsep=2.5mm]
  \item \pp{3} Write down the general form of a \emph{generalised additive model} (GAM) for a
    quantitative response $Y$ and predictors $X_1,\dots,X_p$, and explain in one sentence how it
    generalises multiple linear regression.
  \item \pp{3} State one important \emph{advantage} of GAMs concerning the interpretation of the
    fitted functions $f_j$, and the main structural \emph{limitation} of the additive form ---
    including how this limitation can be repaired.
  \item \pp{2} Describe in one or two sentences how the \emph{backfitting} algorithm fits a GAM.
\end{enumerate}

\ifdefined\withsolutions
\begin{solutionbox}
\textbf{(a) Model form.}
\[
  Y=\beta_0+f_1(X_1)+f_2(X_2)+\dots+f_p(X_p)+\varepsilon .
\]
Multiple linear regression is the special case $f_j(X_j)=\beta_jX_j$; a GAM replaces each linear
component by a smooth non-linear function $f_j$ (e.g.\ a natural or smoothing spline, or local
regression), while keeping the \emph{additive} structure.

\medskip
\textbf{(b) Advantage and limitation.}
\emph{Advantage (interpretability):} because the model is additive, the contribution of each
$X_j$ to $Y$ is captured by its own function $f_j$, which can be examined and plotted
individually while the other variables are held fixed --- non-linear effects remain as easy to
read as in a coefficient table, one panel per predictor.
\emph{Limitation:} additivity excludes \emph{interactions}: the effect of $X_j$ cannot depend
on the value of $X_k$. Repair: add interaction terms manually, e.g.\ a product term
$X_j\times X_k$ or a low-dimensional smooth surface $f_{jk}(X_j,X_k)$ as an extra component.

\medskip
\textbf{(c) Backfitting.}
Cycle through the predictors and repeatedly re-fit each $f_j$ (by a smoother, e.g.\ a spline)
to the \emph{partial residuals} $r_i=y_i-\beta_0-\sum_{k\ne j}f_k(x_{ik})$, holding all other
fitted functions fixed; iterate until the $f_j$ stop changing.
\end{solutionbox}
\fi

% ============================================================
\problem{3: Regression and classification trees by hand}{20}
{\small\itshape Lecture slides: Chapter 8 --- Tree-Based Methods (Lecture 10).}\par\vspace{1mm}

A logistics company records for $n=6$ regional depots the number of delivery vans in operation
$x$ and the number of parcels delivered that day $y$ (in 1{,}000 parcels):

\begin{center}
\begin{tabular}{lcccccc}
\toprule
Depot $i$ & 1 & 2 & 3 & 4 & 5 & 6\\
\midrule
Vans $x_i$ & 1 & 2 & 3 & 4 & 5 & 6\\
Parcels $y_i$ & 3 & 5 & 4 & 9 & 11 & 10\\
\bottomrule
\end{tabular}
\end{center}

A regression tree with a \emph{single} split at cutpoint $s$ partitions the data into
$R_1=\{i: x_i<s\}$ and $R_2=\{i: x_i\ge s\}$ and predicts the region mean $\hat{y}_{R_m}$ in
each region.

\begin{enumerate}[label=(\alph*),leftmargin=7mm,itemsep=2.5mm]
  \item \pp{8} For each of the two candidate cutpoints $s=2.5$ and $s=3.5$, determine the two
    regions, the region means and the resulting total residual sum of squares
    $\operatorname{RSS}(s)=\sum_{m=1}^{2}\sum_{i\in R_m}(y_i-\hat{y}_{R_m})^2$.
    Which cutpoint does recursive binary splitting choose, and why?
  \item \pp{2} Using the chosen split, predict the number of parcels for a new depot that will
    operate $x=5$ vans.
  \item \pp{6} A \emph{classification} tree for a different task has a terminal node containing
    $9$ training observations from three product categories: $5$ ``Electronics'', $3$
    ``Clothing'' and $1$ ``Books''. Compute the class proportions, the Gini index
    $G=\sum_{k}\hat{p}_{k}(1-\hat{p}_{k})$ and the entropy $D=-\sum_{k}\hat{p}_{k}\ln\hat{p}_{k}$
    of this node \emph{(use $\ln(5/9)=-0.5878$, $\ln(1/3)=-1.0986$ and $\ln(1/9)=-2.1972$)}.
    Which class does the node predict?
  \item \pp{4} In \emph{cost-complexity pruning} one minimises, for a subtree
    $T\subseteq T_0$ of the large tree $T_0$,
    \[
      \sum_{m=1}^{|T|}\ \sum_{i:\,x_i\in R_m}\bigl(y_i-\hat{y}_{R_m}\bigr)^2+\alpha\,|T| .
    \]
    Explain the role of the tuning parameter $\alpha$: what happens for $\alpha=0$, what happens
    as $\alpha$ increases, and how is $\alpha$ chosen in practice?
\end{enumerate}

\ifdefined\withsolutions
\begin{solutionbox}
\textbf{(a) Evaluating the two cutpoints.}
Root node for reference: $\bar{y}=42/6=7$.

\emph{Cutpoint $s=2.5$:} $R_1=\{1,2\}$ with $y$-values $\{3,5\}$, mean $\hat{y}_{R_1}=4$;
$R_2=\{3,4,5,6\}$ with $y$-values $\{4,9,11,10\}$, mean $\hat{y}_{R_2}=34/4=8.5$.
\begin{align*}
  \operatorname{RSS}_{R_1}&=(3-4)^2+(5-4)^2=2,\\
  \operatorname{RSS}_{R_2}&=(4-8.5)^2+(9-8.5)^2+(11-8.5)^2+(10-8.5)^2=20.25+0.25+6.25+2.25=29,
\end{align*}
so $\operatorname{RSS}(2.5)=2+29=\mathbf{31}$.

\emph{Cutpoint $s=3.5$:} $R_1=\{1,2,3\}$ with $y$-values $\{3,5,4\}$, mean
$\hat{y}_{R_1}=12/3=4$; $R_2=\{4,5,6\}$ with $y$-values $\{9,11,10\}$, mean
$\hat{y}_{R_2}=30/3=10$.
\[
  \operatorname{RSS}_{R_1}=(3-4)^2+(5-4)^2+(4-4)^2=2,\qquad
  \operatorname{RSS}_{R_2}=(9-10)^2+(11-10)^2+(10-10)^2=2,
\]
so $\operatorname{RSS}(3.5)=2+2=\mathbf{4}$.

Since $4<31$, recursive binary splitting is \emph{greedy} and picks the cutpoint with the
smallest RSS: $\mathbf{s=3.5}$. (It cleanly separates the three low-volume depots, mean $4$,
from the three high-volume depots, mean $10$; in fact $s=3.5$ is the best of all five possible
cutpoints.)

\medskip
\textbf{(b) Prediction for $x=5$.}
$x=5\ge 3.5$, so the new depot falls into $R_2$ and the tree predicts
$\hat{y}=\hat{y}_{R_2}=\mathbf{10}$ (i.e.\ $10{,}000$ parcels).

\medskip
\textbf{(c) Node impurity (three classes).}
$\hat{p}_{\text{E}}=5/9$, $\hat{p}_{\text{C}}=3/9=1/3$, $\hat{p}_{\text{B}}=1/9$.
\[
  G=\tfrac{5}{9}\cdot\tfrac{4}{9}+\tfrac{3}{9}\cdot\tfrac{6}{9}+\tfrac{1}{9}\cdot\tfrac{8}{9}
   =\tfrac{20+18+8}{81}=\tfrac{46}{81}=\mathbf{0.568}.
\]
\[
  D=-\Bigl(\tfrac{5}{9}\ln\tfrac{5}{9}+\tfrac{1}{3}\ln\tfrac{1}{3}+\tfrac{1}{9}\ln\tfrac{1}{9}\Bigr)
   =-\bigl(\tfrac{5}{9}(-0.5878)+\tfrac{1}{3}(-1.0986)+\tfrac{1}{9}(-2.1972)\bigr)
\]
\[
   =0.3266+0.3662+0.2441=\mathbf{0.937}.
\]
The node predicts the majority class, \textbf{Electronics} (with estimated probability
$5/9\approx 0.556$). Both $G$ and $D$ would be $0$ for a pure node; here the node is quite
impure because no class dominates strongly.

\medskip
\textbf{(d) Role of $\alpha$.}
$\alpha$ is the price per terminal node: it trades off the fit of the subtree (first term)
against its complexity $|T|$.
For $\alpha=0$ the criterion equals the training RSS and the full tree $T_0$ is optimal.
As $\alpha$ increases, additional leaves must ``pay for themselves''; branches whose RSS
reduction is smaller than $\alpha$ are pruned, so the optimal subtree becomes smaller ---
one obtains a \emph{nested sequence} of subtrees as $\alpha$ grows (weakest-link pruning).
In practice $\alpha$ is selected by \emph{cross-validation} (choose the $\alpha$ with the
smallest CV error, then refit on all data), which balances bias against variance instead of
optimising training fit.
\end{solutionbox}
\fi

% ============================================================
\problem{4: Ensembles --- bagging, random forests, boosting}{16}
{\small\itshape Lecture slides: Chapter 8 --- Tree-Based Methods (Lecture 10).}\par\vspace{1mm}

\begin{enumerate}[label=(\alph*),leftmargin=7mm,itemsep=2.5mm]
  \item \pp{6} Let $\hat{f}_1(x),\dots,\hat{f}_B(x)$ be $B$ identically distributed (but not
    independent) tree predictions at a fixed point $x$, each with variance $\sigma^2$ and with
    pairwise correlation $\rho$. The variance of the bagged average is
    \[
      \operatorname{Var}\Bigl(\tfrac{1}{B}\sum_{b=1}^{B}\hat{f}_b(x)\Bigr)
      =\rho\,\sigma^{2}+\frac{1-\rho}{B}\,\sigma^{2}.
    \]
    (i) Explain briefly why this formula shows that bagging reduces variance, and which term
    limits the reduction. (ii) Evaluate the variance for $B=50$, $\rho=0.4$ and $\sigma^2=16$,
    and compare it with the variance of a single tree. (iii) What value does the variance
    approach as $B\to\infty$?
  \item \pp{4} How do \emph{random forests} modify bagging at each split, and what is the
    recommended subset size $m$ for \emph{regression} as a function of $p$ (evaluate for
    $p=30$)? Explain, with reference to the formula in (a), why this modification improves on
    bagging.
  \item \pp{4} Name the three tuning parameters of \emph{boosting} for regression trees and
    state in one sentence each what they control. What is the ``slow learning'' idea behind
    boosting?
  \item \pp{2} What does the \emph{out-of-bag} (OOB) error of a bagged model estimate, and why
    is it available essentially for free?
\end{enumerate}

\ifdefined\withsolutions
\begin{solutionbox}
\textbf{(a) Variance of the bagged average.}
(i) The second term $\frac{1-\rho}{B}\sigma^2$ is driven to $0$ by averaging over many trees
($B$ large), so averaging removes the ``independent part'' of the trees' variability ---
this is the variance reduction of bagging. The first term $\rho\sigma^2$ does \emph{not}
depend on $B$: because the bootstrap trees are grown on overlapping samples of the same data
(and share the same strong predictors), they are positively correlated, and this correlated
part of the variance cannot be averaged away. It limits the reduction.
(ii) For $B=50$, $\rho=0.4$, $\sigma^2=16$:
\[
  \operatorname{Var}=0.4\cdot 16+\frac{1-0.4}{50}\cdot 16
  =6.4+\frac{0.6\cdot 16}{50}=6.4+0.192=\mathbf{6.592},
\]
compared with $\sigma^2=16$ for a single tree --- a reduction of about $58.8\%$.
(iii) As $B\to\infty$ the variance approaches the floor $\rho\sigma^2=0.4\cdot 16=\mathbf{6.4}$.
(Increasing $B$ further does not overfit; it only stabilises the average.)

\medskip
\textbf{(b) Random forests.}
At each split, only a \emph{random subset} of $m$ of the $p$ predictors is considered as
splitting candidates; a fresh subset is drawn at every split. Recommended for regression:
$m\approx p/3$, so for $p=30$: $m=30/3=\mathbf{10}$.
Why it helps: with bagging, a dominant predictor is used in the top split of almost every tree,
so the trees look alike and $\rho$ is high. Excluding it from $\approx (p-m)/p$ of the split
searches \emph{decorrelates} the trees, i.e.\ lowers $\rho$ --- and by (a) this lowers the
irreducible floor $\rho\sigma^2$ of the ensemble variance, which mere averaging cannot touch.

\medskip
\textbf{(c) Boosting.}
\begin{itemize}[nosep,leftmargin=5mm]
  \item \emph{Number of trees $B$:} how many small trees are added sequentially; unlike
    bagging, boosting \emph{can} overfit if $B$ is too large, so $B$ is chosen by
    cross-validation.
  \item \emph{Shrinkage $\lambda$} (learning rate, e.g.\ $0.01$ or $0.001$): scales the
    contribution of each new tree and thus the speed of learning; smaller $\lambda$ usually
    needs larger $B$.
  \item \emph{Number of splits $d$ per tree} (interaction depth): the complexity of each tree
    and the interaction order of the model; often even $d=1$ (stumps, additive model) works
    well.
\end{itemize}
\emph{Slow learning:} each tree is fitted to the current \emph{residuals}, and only a small
fraction $\lambda$ of it is added to the model; the model thus improves slowly and exactly in
the regions where it still performs poorly, instead of fitting the data hard in one step ---
which tends to generalise better.

\medskip
\textbf{(d) OOB error.}
Each bootstrap sample leaves out on average about $1/3$ of the observations
(``out-of-bag''). Predicting each observation only with the trees that did \emph{not} use it,
and aggregating these predictions, yields the OOB error --- a valid estimate of the
\emph{test error} of the bagged model. It is free because it reuses the bootstrap side
product; no separate validation set or cross-validation loop is required.
\end{solutionbox}
\fi

% ============================================================
\problem{5: Neural networks by hand}{20}
{\small\itshape Lecture slides: Chapter 10 --- Deep Learning (Lecture 11).}\par\vspace{1mm}

Consider a tiny feed-forward network for a quantitative response with input
$x=(x_1,x_2)=(2,3)$, one hidden layer with two ReLU units ($g(z)=\max\{0,z\}$) and a linear
output unit:
\[
  A_1=g\bigl(-3+2\,x_1+1\,x_2\bigr),\qquad
  A_2=g\bigl(2-2\,x_1-1\,x_2\bigr),\qquad
  f(x)=-2+3\,A_1+2\,A_2 .
\]

\begin{enumerate}[label=(\alph*),leftmargin=7mm,itemsep=2.5mm]
  \item \pp{8} Perform the forward pass step by step: compute both pre-activations, both
    activations $A_1,A_2$ and the network output $f(x)$. Comment briefly on what the ReLU did
    to the second unit and why such non-linear activations are essential.
  \item \pp{6} A classification network with $K=3$ classes produces the logits (output-layer
    pre-activations) $z=(z_1,z_2,z_3)=(2.0,\ 0.0,\ -1.0)$ for one observation whose true class
    is class~2. Compute all three softmax probabilities and the cross-entropy loss
    $-\ln \hat{p}_{\text{true}}$.
    \emph{(Use $e^{2}=7.389$, $e^{0}=1.000$, $e^{-1}=0.368$ and $\ln 8.757=2.170$.)}
  \item \pp{4} Count the parameters of a fully connected network with $p=6$ inputs, one hidden
    layer with $k=5$ units and an output layer with $K=3$ units, where every hidden and output
    unit has a bias. Show the count separately for the two weight layers.
  \item \pp{2} A convolutional layer receives a $64\times 64$ image and uses filters of size
    $F=3$ with padding $P=1$ and stride $S=2$. Compute the spatial size of the output feature
    map from $\bigl\lfloor(W-F+2P)/S\bigr\rfloor+1$ (note the stride is $2$, so the division is
    rounded down).
\end{enumerate}

\ifdefined\withsolutions
\begin{solutionbox}
\textbf{(a) Forward pass.}
Pre-activations:
\[
  z_1=-3+2\cdot 2+1\cdot 3=-3+4+3=4,\qquad
  z_2=2-2\cdot 2-1\cdot 3=2-4-3=-5 .
\]
ReLU activations:
\[
  A_1=g(4)=\max\{0,4\}=4,\qquad A_2=g(-5)=\max\{0,-5\}=0 .
\]
Output:
\[
  f(x)=-2+3\cdot 4+2\cdot 0=-2+12+0=\mathbf{10}.
\]
The ReLU \emph{switched the second unit off}: its negative pre-activation is clipped to $0$,
so unit 2 contributes nothing for this input (but would for other inputs). This thresholding
is the source of non-linearity: without a non-linear $g$, the network would collapse into a
single linear function of $x$, no matter how many layers it has.

\medskip
\textbf{(b) Softmax and cross-entropy.}
Exponentials: $e^{2.0}=7.389$, $e^{0.0}=1.000$, $e^{-1.0}=0.368$; their sum is
$7.389+1.000+0.368=8.757$. Softmax probabilities:
\[
  \hat{p}_1=\frac{7.389}{8.757}=0.844,\qquad
  \hat{p}_2=\frac{1.000}{8.757}=0.114,\qquad
  \hat{p}_3=\frac{0.368}{8.757}=0.042,
\]
which sum to $1.000$. Cross-entropy loss for the true class $2$:
\[
  -\ln\hat{p}_2=-\ln\frac{1.000}{8.757}=-\bigl(0-\ln 8.757\bigr)=\ln 8.757=\mathbf{2.170}.
\]
(Note that the model's most probable class is class~1, but the loss is evaluated at the
\emph{true} class~2, which received only probability $0.114$ --- hence the sizeable loss.)

\medskip
\textbf{(c) Parameter count for $6\to 5\to 3$.}
Input $\to$ hidden: $6\cdot 5=30$ weights $+\,5$ biases $=35$.
Hidden $\to$ output: $5\cdot 3=15$ weights $+\,3$ biases $=18$.
Total: $35+18=\mathbf{53}$ parameters.

\medskip
\textbf{(d) Convolution output size.}
\[
  \Bigl\lfloor\frac{W-F+2P}{S}\Bigr\rfloor+1
  =\Bigl\lfloor\frac{64-3+2\cdot 1}{2}\Bigr\rfloor+1
  =\Bigl\lfloor\frac{63}{2}\Bigr\rfloor+1=\lfloor 31.5\rfloor+1=31+1=32,
\]
so the output feature map is $\mathbf{32\times 32}$ --- the stride $S=2$ roughly halves the
$64\times 64$ input (the floor discards the fractional part before adding $1$).
\end{solutionbox}
\fi

% ============================================================
\problem{6: Multiple testing}{20}
{\small\itshape Lecture slides: Chapter 13 --- Multiple Testing (Lecture 12).}\par\vspace{1mm}

\begin{enumerate}[label=(\alph*),leftmargin=7mm,itemsep=2.5mm]
  \item \pp{4} A lab performs $m=40$ \emph{independent} hypothesis tests, each at level
    $\alpha=0.05$, and suppose all $40$ null hypotheses are true. Define the family-wise error
    rate (FWER) and compute it for this situation \emph{(use $(0.95)^{40}=0.1285$)}. Interpret
    the result in one sentence.
  \item \pp{7} A research group screens $m=6$ candidate biomarkers; the ordered $p$-values are
    \begin{center}
    \begin{tabular}{lcccccc}
    \toprule
    $j$ & 1 & 2 & 3 & 4 & 5 & 6\\
    \midrule
    $p_{(j)}$ & 0.001 & 0.008 & 0.011 & 0.030 & 0.041 & 0.260\\
    \bottomrule
    \end{tabular}
    \end{center}
    Control the FWER at $\alpha=0.05$ (i) with \emph{Bonferroni} and (ii) with the
    \emph{Holm step-down} procedure. Give the relevant thresholds, state for each method
    exactly which hypotheses are rejected, and explain why Holm can never reject fewer
    hypotheses than Bonferroni.
  \item \pp{6} Apply the \emph{Benjamini--Hochberg} (BH) procedure with FDR level $q=0.05$ to
    the same six $p$-values: give the thresholds $jq/m$, determine which hypotheses are
    rejected, and state precisely what ``FDR control at $q=0.05$'' guarantees for the
    resulting set of rejections.
  \item \pp{3} A pharmaceutical company runs a \emph{confirmatory} phase-III trial that tests
    $m=5$ pre-specified endpoints; a single false-positive endpoint could lead to the approval
    of an ineffective drug. Should they control the FWER or the FDR? Justify your
    recommendation.
\end{enumerate}

\ifdefined\withsolutions
\begin{solutionbox}
\textbf{(a) FWER for $m=40$ independent tests.}
$\text{FWER}=\Pr(\text{at least one Type I error})=\Pr(V\ge 1)$, where $V$ is the number of
true nulls that are (falsely) rejected. With independence and all nulls true:
\[
  \text{FWER}=1-\Pr(\text{no false rejection})=1-(1-0.05)^{40}=1-(0.95)^{40}
  =1-0.1285=\mathbf{0.8715}.
\]
Interpretation: even though each single test is performed at the $5\%$ level, there is about an
$87\%$ chance of at least one false discovery among the $40$ tests --- testing many hypotheses
without adjustment almost guarantees spurious findings.

\medskip
\textbf{(b) Bonferroni and Holm at $\alpha=0.05$.}
\emph{Bonferroni:} reject $H_{(j)}$ iff $p_{(j)}\le\alpha/m=0.05/6=0.00833$.
$p_{(1)}=0.001\le 0.00833$ \checkmark, $p_{(2)}=0.008\le 0.00833$ \checkmark,
$p_{(3)}=0.011>0.00833$ $\times$.
$\Rightarrow$ Bonferroni rejects $H_{(1)}$ and $H_{(2)}$: \textbf{2 rejections}.

\emph{Holm:} compare $p_{(j)}$ with $\alpha/(m-j+1)$ step by step and stop at the first failure:
\begin{center}
\begin{tabular}{lcccccc}
\toprule
$j$ & 1 & 2 & 3 & 4 & 5 & 6\\
\midrule
$p_{(j)}$ & 0.001 & 0.008 & 0.011 & 0.030 & 0.041 & 0.260\\
$\alpha/(m-j+1)$ & 0.00833 & 0.0100 & 0.0125 & 0.0167 & 0.0250 & 0.0500\\
reject? & \checkmark & \checkmark & \checkmark & stop & --- & ---\\
\bottomrule
\end{tabular}
\end{center}
$0.001\le 0.00833$, $0.008\le 0.0100$, $0.011\le 0.0125$, but $0.030>0.0167$: stop.
$\Rightarrow$ Holm rejects $H_{(1)},H_{(2)},H_{(3)}$: \textbf{3 rejections}.

\emph{Why Holm $\supseteq$ Bonferroni:} Holm's thresholds $\alpha/(m-j+1)$ are $\ge\alpha/m$
for every $j$ (with equality only at $j=1$), so every $p$-value that passes Bonferroni also
passes Holm's less stringent later hurdles --- Holm is uniformly more powerful, yet still
controls the FWER at $\alpha$ without any independence assumption. Here Holm rejects one more
hypothesis ($H_{(3)}$) than Bonferroni.

\medskip
\textbf{(c) Benjamini--Hochberg at $q=0.05$.}
Thresholds $jq/m=j\cdot 0.05/6$:
\begin{center}
\begin{tabular}{lcccccc}
\toprule
$j$ & 1 & 2 & 3 & 4 & 5 & 6\\
\midrule
$p_{(j)}$ & 0.001 & 0.008 & 0.011 & 0.030 & 0.041 & 0.260\\
$jq/m$ & 0.00833 & 0.0167 & 0.0250 & 0.0333 & 0.0417 & 0.0500\\
$p_{(j)}\le jq/m$? & yes & yes & yes & yes & \textbf{yes} & no\\
\bottomrule
\end{tabular}
\end{center}
$L=$ largest $j$ with $p_{(j)}\le jq/m$: $j=5$ (since $0.041\le 0.0417$, while
$0.260>0.0500$). BH rejects \emph{all} $H_{(j)}$ with $j\le L$: $H_{(1)},\dots,H_{(5)}$:
\textbf{5 rejections}.
Guarantee: $\text{FDR}=E\bigl[V/\max(R,1)\bigr]\le q=0.05$ --- \emph{on average} at most
$5\%$ of the rejected hypotheses are false discoveries. It does \emph{not} guarantee that at
most $5\%$ of these particular $5$ rejections are false, nor does it bound the probability of
making any false rejection (that would be FWER control). Note BH is the most liberal of the
three here ($5$ rejections vs.\ $3$ for Holm and $2$ for Bonferroni).

\medskip
\textbf{(d) Confirmatory phase-III trial.}
Control the \textbf{FWER}. This is a confirmatory study with only $m=5$ pre-specified
endpoints, and the cost of even a single false positive is severe (approval of an ineffective
drug, patient harm, regulatory failure). FWER control (e.g.\ Holm at $\alpha=0.05$) bounds the
probability of \emph{any} false rejection at $5\%$, which is exactly the guarantee a regulator
demands. FDR control would only limit the \emph{expected proportion} of false endpoints and
would tolerate an occasional false positive --- acceptable in large exploratory screens, but
not when each individual claim must stand on its own.
\end{solutionbox}
\fi

% ============================================================
\problem{7: Cumulative essentials (Chapters 2--6)}{20}
{\small\itshape Lecture slides: Chapters 2--6 (Lectures 1--8).}\par\vspace{1mm}

\begin{enumerate}[label=(\alph*),leftmargin=7mm,itemsep=2.5mm]
  \item \pp{4} The \emph{bootstrap} draws a resample of size $n$ with replacement from the $n$
    observations. (i) Show that the probability that a given observation is \emph{not} selected
    into a particular bootstrap sample equals $(1-1/n)^n$, and give its limit as $n\to\infty$
    \emph{(use $e^{-1}=0.368$)}. (ii) Hence what fraction of the original observations appears
    (at least once) in a typical bootstrap sample as $n\to\infty$? (iii) Compute this inclusion
    probability exactly for $n=5$ \emph{(use $0.8^{5}=0.328$)}.
  \item \pp{4} At a fixed test point $x_0$ the expected test error of an estimate $\hat f$
    decomposes as
    $E\bigl[(y_0-\hat f(x_0))^2\bigr]
      =\operatorname{Var}(\hat f(x_0))+\bigl[\operatorname{Bias}(\hat f(x_0))\bigr]^2
       +\operatorname{Var}(\varepsilon)$.
    Two candidate methods give, at $x_0$:
    \begin{center}
    \begin{tabular}{lccc}
    \toprule
    Method & $\operatorname{Var}(\hat f)$ & $\operatorname{Bias}(\hat f)$ & $\operatorname{Var}(\varepsilon)$\\
    \midrule
    (1) flexible spline & 3.5 & $-1.0$ & 2.0\\
    (2) linear model & 0.5 & $-2.5$ & 2.0\\
    \bottomrule
    \end{tabular}
    \end{center}
    Compute the expected test MSE of each method, say which you would deploy at $x_0$, and state
    the lowest test MSE \emph{any} method could achieve here.
  \item \pp{3} Write down the \emph{elastic net} penalty and explain what it inherits from the
    lasso and what it inherits from ridge regression. In which situation does the elastic net
    clearly outperform the lasso?
  \item \pp{3} Two predictors in a regression have correlation $0.95$. Explain what such
    \emph{collinearity} does to the least-squares coefficient estimates. The variance inflation
    factor is $\mathrm{VIF}_j=1/(1-R_j^2)$, where $R_j^2$ is from regressing $X_j$ on the other
    predictors; compute and interpret $\mathrm{VIF}_j$ for $R_j^2=0.9$ \emph{(use
    $\sqrt{10}=3.162$)}.
  \item \pp{6} Match each scenario to exactly one method from
    \{\emph{linear regression}, \emph{lasso}, \emph{logistic regression},
    \emph{random forest}\} (each method used exactly once) and justify each choice in one
    sentence:
    \begin{enumerate}[label=(\roman*),nosep,topsep=1mm]
      \item Estimate, with a $95\%$ confidence interval, how much one additional year of
        schooling raises hourly wage, from $n=3{,}000$ workers and a handful of predictors.
      \item From $n=80$ patients with $p=8{,}000$ metabolite measurements, predict a continuous
        biomarker level and identify a short list of relevant metabolites.
      \item Predict whether a customer will churn (yes/no) and report to management how each
        customer characteristic changes the \emph{odds} of churning.
      \item Forecast next-day electricity demand (continuous) from hundreds of weather and
        calendar variables with strong non-linear effects and interactions; only forecast
        accuracy matters.
    \end{enumerate}
\end{enumerate}

\ifdefined\withsolutions
\begin{solutionbox}
\textbf{(a) Bootstrap inclusion probability.}
(i) In each of the $n$ independent draws, a fixed observation is \emph{missed} with probability
$1-\tfrac1n$; since the draws are independent, all $n$ draws miss it with probability
$(1-\tfrac1n)^n$. As $n\to\infty$, $(1-\tfrac1n)^n\to e^{-1}=\mathbf{0.368}$.
(ii) The observation is therefore \emph{included} with probability $1-e^{-1}=1-0.368=0.632$, so
on average about $\mathbf{63.2\%}$ of the original observations appear in a bootstrap sample
(the remaining $\approx 1/3$ are the ``out-of-bag'' observations).
(iii) For $n=5$: $\Pr(\text{included})=1-(1-\tfrac15)^5=1-0.8^{5}=1-0.328=\mathbf{0.672}$.

\medskip
\textbf{(b) Bias--variance decomposition.}
\[
  \text{MSE}_1=\operatorname{Var}+\operatorname{Bias}^2+\operatorname{Var}(\varepsilon)
   =3.5+(-1.0)^2+2.0=3.5+1.0+2.0=\mathbf{6.5},
\]
\[
  \text{MSE}_2=0.5+(-2.5)^2+2.0=0.5+6.25+2.0=\mathbf{8.75}.
\]
Deploy the \textbf{flexible spline} (method 1): its higher estimation variance ($3.5$ vs.\
$0.5$) is more than offset by its much smaller squared bias ($1.0$ vs.\ $6.25$), giving the
lower expected test error ($6.5<8.75$). The lowest test MSE any method can achieve at $x_0$ is
the \emph{irreducible error} $\operatorname{Var}(\varepsilon)=\mathbf{2.0}$; no estimator can go
below it.

\medskip
\textbf{(c) Elastic net.}
The elastic net minimises, for $0\le\alpha\le 1$,
\[
  \sum_{i=1}^n\bigl(y_i-\beta_0-\textstyle\sum_j\beta_jx_{ij}\bigr)^2
  +\lambda\Bigl[\tfrac{1-\alpha}{2}\textstyle\sum_j\beta_j^2+\alpha\textstyle\sum_j|\beta_j|\Bigr],
\]
i.e.\ a weighted mix of the ridge ($\ell_2$) and lasso ($\ell_1$) penalties. From the
\emph{lasso} it inherits the ability to set coefficients \emph{exactly to zero} (sparsity /
variable selection); from \emph{ridge} it inherits the graceful handling of correlated
predictors and the ability to select more than $n$ variables when $p>n$. It clearly beats the
lasso when there are \emph{groups of highly correlated predictors}: the lasso tends to pick one
member of a group essentially at random and drop the rest, whereas the elastic net's ridge part
shares the weight across the group and keeps them in or out together (the ``grouping effect'').

\medskip
\textbf{(d) Collinearity and VIF.}
Collinearity does not bias the OLS estimates, but it makes them \emph{imprecise and unstable}:
it inflates their variances and standard errors, widens confidence intervals, and can make
individual coefficients change sign or lose significance even when the predictors are jointly
important (predictions may still be fine). With $R_j^2=0.9$,
\[
  \mathrm{VIF}_j=\frac{1}{1-0.9}=\frac{1}{0.1}=\mathbf{10},
\]
so the variance of $\hat\beta_j$ is $10$ times larger than it would be with an uncorrelated
$X_j$, and its standard error is inflated by $\sqrt{10}=\mathbf{3.162}$. As a rule of thumb a
$\mathrm{VIF}$ above $5$--$10$ signals problematic collinearity.

\medskip
\textbf{(e) Matching.}
\begin{itemize}[nosep,leftmargin=5mm]
  \item (i) \textbf{Linear regression}: quantitative response, goal is \emph{inference}
    (effect size with confidence interval), $n\gg p$ --- exactly what OLS with standard errors
    provides.
  \item (ii) \textbf{Lasso}: $p\gg n$ ($8{,}000$ vs.\ $80$) makes OLS infeasible; the $\ell_1$
    penalty regularises and selects a short list of metabolites.
  \item (iii) \textbf{Logistic regression}: binary response with interpretable coefficients
    ($e^{\hat\beta_j}$ = odds ratios) --- suitable for reporting how each characteristic changes
    the odds.
  \item (iv) \textbf{Random forest}: captures non-linearities and interactions automatically
    from many predictors and is accurate and robust; interpretability is not required, only
    forecast accuracy.
\end{itemize}
\end{solutionbox}
\fi

\vspace{6mm}
\begin{center}
\rule{0.35\textwidth}{0.4pt}\\[1mm]
\textit{End of the examination --- good luck!}
\end{center}

\end{document}
